% Standalone source assembled by build.py.

% BEGIN preamble.tex
\documentclass[11pt,letterpaper]{article}
\usepackage[T1]{fontenc}
\usepackage[utf8]{inputenc}
\usepackage{lmodern}
\usepackage[margin=1in,headheight=15pt]{geometry}
\usepackage{amsmath,amssymb,amsthm,mathtools,mathrsfs}
\usepackage{microtype}
\usepackage{booktabs,longtable,array}
\usepackage{xcolor}
\usepackage{enumitem}
\usepackage{fancyhdr}
\usepackage{hyperref}
\usepackage{bookmark}
\definecolor{inkblue}{HTML}{173A56}
\definecolor{linkblue}{HTML}{165F76}
\hypersetup{colorlinks=true,linkcolor=inkblue,citecolor=linkblue,
 urlcolor=linkblue,pdftitle={Lattice Contacts and Mixed Zeta Identities},
 pdfauthor={Research continuation prepared for the ProveIt project},
 pdfsubject={Exact identities for Hurwitz tails, Stieltjes products, and Gamma spectral faces}}
\urlstyle{same}
\setlength{\emergencystretch}{2em}
\setlength{\parskip}{3pt}
\setlist{topsep=4pt,itemsep=4pt}
\numberwithin{equation}{section}
\newtheorem{theorem}{Theorem}[section]
\newtheorem{lemma}[theorem]{Lemma}
\newtheorem{proposition}[theorem]{Proposition}
\newtheorem{corollary}[theorem]{Corollary}
\theoremstyle{definition}
\newtheorem{definition}[theorem]{Definition}
\newtheorem{conjecture}[theorem]{Conjecture}
\theoremstyle{remark}
\newtheorem{remark}[theorem]{Remark}
\DeclareMathOperator{\Li}{Li}
\DeclareMathOperator{\Log}{Log}
\DeclareMathOperator{\FP}{FP}
\DeclareMathOperator{\Ree}{Re}
\newcommand{\ii}{\mathrm i}
\newcommand{\dd}{\,\mathrm d}
\newcommand{\C}{\mathbb C}
\newcommand{\N}{\mathbb N}
\newcommand{\Q}{\mathbb Q}
\newcommand{\R}{\mathbb R}
\pagestyle{fancy}
\fancyhf{}
\fancyhead[L]{\small\textsc{Lattice Contacts and Mixed Zeta Identities}}
\fancyhead[R]{\small ProveIt research}
\fancyfoot[C]{\thepage}
\renewcommand{\headrulewidth}{0.3pt}
\allowdisplaybreaks[2]


% END INPUT

\begin{document}
\hypersetup{pageanchor=false}
\begin{titlepage}
\centering
\vspace*{0.7in}
{\large\color{linkblue}\textsc{ProveIt research continuation}\par}
\vspace{0.35in}
{\Huge\bfseries\color{inkblue} Lattice Contacts and\\[5pt] Mixed Zeta Identities\par}
\vspace{0.22in}
{\Large Unequal harmonic tails, shifted Stieltjes products,\\[3pt]
and Gamma spectral faces\par}
\vspace{0.4in}
{\large 11 October 2026\par}
\vspace{0.2in}
{\normalsize Prepared for the ProveIt project\par}
\vspace{0.4in}
\begin{minipage}{0.9\textwidth}
\textbf{Abstract.}
We continue several explicit research questions in the ProveIt
manuscript and its latest incoming reports. A single elementary
generating function gives every contact term at minimally compensated
nonpositive even orders of a vertical Hurwitz lattice; periodic
weights remove all these terms exactly when their mean is zero.
An order-dependent Mellin identity supplies exact modular transport
and every nonmodular derivative correction. A polarized Bernoulli
identity extends the centered-moment calculus to arbitrarily many
unequal even-order Hurwitz tails, with terminating recurrences,
finite multiple-zeta spans, and explicit six- and eight-coordinate
examples. Independently shifted products of generalized Stieltjes
constants are evaluated by a jointly holomorphic generator with
geometrically convergent Hurwitz series. Finally, every polynomial
spectral-face coefficient is explicit, every Gamma face with one
displaced mass reduces to depth-two polylogarithms, and the precise
polar obstruction for independently variable spectral orders is
classified and removed. All claimed identities have analytic or
finite algebraic proofs. Numerical checks serve as independent
diagnostics. A source-status correction, integration guidance, and
a prioritized research agenda are included.
\end{minipage}
\vfill
{\small Source snapshot: \texttt{36edefa8c531d9d6df8306ac51bcc317250f404c}\par}
{\small Novelty is assessed against the inspected corpus and specified
primary literature; historical priority is not asserted.\par}
\end{titlepage}
\hypersetup{pageanchor=true}
\pagenumbering{roman}
\tableofcontents
\clearpage
\pagenumbering{arabic}

% BEGIN sections/01_overview.tex
\section{Research frontier, principal identities, and conventions}
\label{sec:overview}

This article develops exact identities at the meeting point of
Hurwitz continuation, polynomially weighted harmonic tails, Gamma
products, and polylogarithms. The source is the ProveIt repository
at commit
\texttt{36edefa8c531d9d6df8306ac51bcc317250f404c}.
The manuscript and the ten incoming archives present at that snapshot
were used to identify the current frontier
\cite{RepoManuscript,RepoIncoming}. The proofs below are self-contained
apart from standard special-function identities and explicitly
credited predecessor formulas that are restated before use.

Two short consequences convey the character of the results. Put
\[
 x_n=n+\tfrac12,\qquad
 T_p(n)=\zeta(p,n+1)=\zeta(p)-H_n^{(p)}.
\]
Then
\begin{equation}\label{eq:overview-cubic}
 \boxed{\sum_{n=0}^{\infty}T_2(n)T_4(n)
                  \bigl(1-x_nT_2(n)\bigr)=\frac5{16}\zeta(6),}
\end{equation}
and
\begin{equation}\label{eq:overview-quartic}
 \boxed{\begin{aligned}
 &\sum_{n=0}^{\infty}T_2(n)^2T_4(n)
       \left\{\left(x_n^2+\frac1{12}\right)T_2(n)-x_n\right\}\\
 &\hspace{13mm}
 =\frac{59}{12}\zeta(3)\zeta(4)
       +\frac{23}{6}\zeta(2)\zeta(5)-\frac{203}{16}\zeta(7).
 \end{aligned}}
\end{equation}
Both sums are ordinary and absolutely convergent. Their summands
decay respectively as \(O(x_n^{-6})\) and \(O(x_n^{-8})\).
They follow from the same finite recurrence that computes all
regular centered moments of any product of positive even-order tails;
see Corollaries~\ref{mtp:short-cubic} and \ref{mtp:short-quartic}.
No finite-part convention is involved in these two formulas.

\subsection{What is proved}

\begin{enumerate}[label=\textbf{\arabic*.},leftmargin=*]
\item \textbf{Every minimal lattice contact.}
At \(s=-N\), \(N\) even, the \(d\)-th order derivative of the
minimally compensated Hurwitz lattice has a complete polynomial
in \(\log(hK)\), where \(K\) is the integer cutoff. Its generating
function is
\[
 \frac{\tau\beta_N}{2}\,
 u\prod_{j=1}^{N}(1-u/j)
 \frac{\sin(\pi u/2)}{\pi u/2}e^{-Xu},
 \qquad
 \tau=\frac{2\pi}{h},\quad
 \beta_N=\frac{B_{N+2}}{(N+1)(N+2)}.
\]
The coefficient multiplied by \(d!\) is the contact of order \(d\).
This answers the all-derivative and filter questions in the
vertical-lattice report \cite{RepoVertical}.

\item \textbf{Exact order-dependent modular transport.}
The Lambert family \(\Lambda(s,\tau)\) in
\eqref{eq:mod-lambda}
has a holomorphic contour transformation near every
nonpositive even order and near \(s=1\).
An explicit trigonometric polynomial recurrence produces all
derivative corrections. The zeroth-order eta and Ramanujan
transformations are credited classical specializations;
the higher derivatives retain specified convergent contour terms.

\item \textbf{Arbitrarily many unequal even tails.}
For a finite labeled collection of even integers \(p_i\ge2\),
every regular centered moment of
\(\prod_iT_{p_i}\) is computed by a finite recurrence on proper
subsets. Its generating function is analytic for \(|t|<2\pi\).
All moments lie in a fixed rational span of convergent multiple
zeta values of weight at most
\[
 \sum_i p_i-1-\mathbf1_{\{|I|\ {\rm odd}\}}.
\]
The mixed products \(T_2^2T_4\) and \(T_2^3T_4\) have explicit
generators and full rank on six and eight specified formal
coordinates. Their first six odd and first eight even moments,
respectively, reconstruct every later regular moment.

\item \textbf{Independent nonperiodic Stieltjes shifts.}
For \(\Ree a,\Ree b>0\), the integral
\[
 \int_0^1
   \left(\zeta(1+u,a+x)-\frac1u\right)
   \left(\zeta(1+v,b+x)-\frac1v\right)\dd x
\]
is evaluated by a finite incomplete-beta prefix, explicitly
removable Gamma quotients, and two geometric Hurwitz series.
The expression is jointly holomorphic in both orders.
Every product of generalized Stieltjes constants and their
argument derivatives follows by coefficient extraction.
Ordinary polygamma products admit a stronger terminating
reduction to two logarithmic Hurwitz sums.

\item \textbf{All spectral-face coefficients and a depth-two collapse.}
A finite formula in harmonic star sums gives
\(I_k(B\mapsto B^m;b_1,\ldots,b_k)\) for every \(k,m\).
With the exact face normalization of
Section~\ref{sfj:sec:faces}, one has
\[
 \begin{aligned}
 &I_k\bigl(G_a;\underbrace{B,\ldots,B}_{k-1},B+\delta\bigr)\\
 &\quad=(-1)^{k-1}(k-1)!
 \sum_{n=0}^{\infty}
 \Li_{1,k-1}\left(\frac{\delta}{(n+a)^2+B+\delta}\right),
 \end{aligned}
\]
where \(k\ge2\), \(a>0\), \(B,\delta\ge0\), and
\(G_a(B)=\log\Gamma(a+\ii\sqrt B)+
\log\Gamma(a-\ii\sqrt B)-\log(2\pi)\).
Thus one displaced mass has functional depth two in every
simplex dimension, with no small-displacement restriction.

\item \textbf{The joint-order obstruction and its exact correction.}
For polynomially weighted products of quadratic spectra, the only
possible polar divisor through the origin is
\(\sum_i\alpha_i=0\). Its numerator is an explicit finite
polynomial. For unequal masses and no periodic multiplier, the joint germ is
holomorphic exactly when the polynomial weight \(W(x)\) is even.
A mean-zero periodic multiplier removes the obstruction for
every polynomial. After the specified local subtraction, all
mixed derivatives have convergent simplex formulas.
\end{enumerate}

\subsection{Relation to established work}

The unweighted product sums of zeta tails already have finite
multiple-zeta evaluations; Hoffman gives a general treatment
\cite{Hoffman}. Later product-sum work includes Li and Chu
\cite{LiChu}. Our new tail assertion concerns the complete
centered polynomial-moment system for arbitrary unequal even
orders, including its subtraction constants and anchored analytic
generator. The displayed short identities are useful consequences
of that larger system.

The incomplete-beta method for products of Hurwitz zeta functions
has a direct predecessor in Shpot and Paris
\cite{ShpotParis}, and logarithmic digamma-square integrals
were studied by Connon \cite{Connon}. The additional content
here is the simultaneous treatment of two independent ordinary
argument shifts, joint removal of both order poles, and an
explicit derivative calculus with proved normal convergence.
Likewise, Dirichlet averages and Gamma products are classical
\cite{Carlson,DLMF}. They underlie the new
collision coefficients and displaced-mass evaluation rather than
constitute novelty claims.

Vertical Gamma-product identities also predate this continuation,
including recent work of Arabi Ardehali and Rosengren
\cite{ArdehaliRosengren}. Generalized Stieltjes modular relations
are established in the work of Dixit, Sathyanarayana, and Sharan
\cite{DixitStieltjes}. The present lattice and modular sections
specify the all-order boundary polynomials, their cutoff
coordinates, and the exact contour corrections for this particular
sampling family. The literature comparison is targeted. It
supports a continuation of the inspected ProveIt questions, not
an exhaustive assertion of historical priority.

\subsection{Conventions and the logical status of certificates}

We write \(\zeta^{(r)}(s,a)\) for differentiation in the order
\(s\), and \(\psi^{(r)}(a)\) for differentiation in the argument.
Argument derivatives of generalized Stieltjes constants are
explicitly indicated. The convention is
\[
 \zeta(1+u,a)=\frac1u+
       \sum_{m\ge0}\frac{(-1)^m}{m!}\gamma_m(a)u^m,
 \qquad \gamma_0(a)=-\psi(a).
\]
We use \(B_1=-1/2\) and
\(H_N^{(r)}=\sum_{j=1}^N j^{-r}\), with \(H_0^{(r)}=0\).
Multiple zeta values have decreasing summation indices;
\(\Li_{s_1,\ldots,s_\ell}(z)\) places \(z^{n_1}\) on the
largest index. Powers and logarithms use the analytic branch
specified on each domain.

An ordinary sum, an analytic continuation value, and a constant
term in a cutoff expansion are identified only after a proof
relates them. Changing the cutoff coordinate can change a
finite part. Similarly, a germ holomorphic along common rays
need not be jointly holomorphic in independent variables.
These distinctions explain the contact polynomials and the
multivariate counterexamples below.

Every theorem is proved analytically or by finite exact algebra.
Rational determinant certificates establish rank on declared
formal coordinates; they do not assert arithmetic independence
of zeta values. High-precision calculations independently check
the formulas, including complex shifts and degenerate limits,
but are not interval certificates or substitutes for proof.
The Gaussian \(S_6\) and revised \(S_8\) identities in the
manuscript remain conjectural \cite{RepoManuscript}.


% END INPUT


% BEGIN sections/02_lattice_contacts.tex
\section{Every minimally compensated vertical-lattice derivative}
\label{sec:lattice}

The vertical-lattice report \cite{RepoVertical} proves a normally
convergent sampling formula after Euler--Maclaurin compensation, and
computes a single extra Bernoulli term at the first negative-even
derivative. Its Questions 1 and 6 ask for all higher minimally
compensated derivatives and for the periodic filters that remove their
contacts. We answer both questions. The answer is a single polynomial
in the logarithm of the cutoff, specified by an elementary generating
function. In particular, higher derivative order does not require
additional moments of a periodic filter.

\subsection{The inherited sampling identity and its exact domain}

Let \(h>0\), \(\tau=2\pi/h\), and \(\Ree a>0\). All powers use the
analytic logarithm on the right half-plane. Write
\[
 w_n=a+\ii hn,\qquad q=e^{-\tau},\qquad z=e^{-\tau a}.
\]
The generalized Stieltjes constants have the convention
\begin{equation}\label{eq:lattice-stieltjes}
 \zeta(1+u,w)=\frac1u+
       \sum_{m=0}^{\infty}\frac{(-1)^m}{m!}\gamma_m(w)u^m,
 \qquad \gamma_0(w)=-\psi(w).
\end{equation}
We use \(B_1=-1/2\). For \(M\geq0\), define
\begin{align}
 R_M(s,w)&=\zeta(s,w)-\frac{w^{1-s}}{s-1}-\frac12w^{-s}
 -\sum_{j=1}^{M}\frac{B_{2j}}{(2j)!}(s)_{2j-1}w^{1-s-2j},
 \label{eq:root-remainder}\\
 {\cal L}_{\nu}(q,z)&=\sum_{k=1}^{\infty}
             \frac{z^k}{k^\nu(1-q^k)},\label{eq:root-lambert}\\
 {\cal B}_M(s;a,\tau)&={\cal L}_{1-s}(q,z)
       -\frac12\Li_{1-s}(z)-\frac1\tau\Li_{2-s}(z)\notag\\
 &\hspace{7mm}
       -\sum_{j=1}^{M}\frac{B_{2j}\tau^{2j-1}}{(2j)!}
                    \Li_{2-s-2j}(z),\label{eq:root-bracket}\\
 {\cal F}_M(s;a,h)&=\frac{\tau^s}{\Gamma(s)}{\cal B}_M(s;a,\tau).
 \label{eq:root-spectral}
\end{align}
The removable value at \(s=1\) is understood in \(R_M\).
Both \(R_M(s,w)\) and \({\cal F}_M(s;a,h)\) are entire in \(s\).

For completeness, the previously established sampling formula is
\begin{equation}\label{eq:root-master}
 \sum_{n\in\mathbb Z}R_M(s,a+\ii hn)={\cal F}_M(s;a,h),
 \qquad \Ree s>-2M.
\end{equation}
It includes all fixed order derivatives on compact subsets of that
strict half-plane. One direct proof applies Poisson periodization to
\[
 \frac{t^{s-1}e^{-at}}{\Gamma(s)}
 \left(\frac1{1-e^{-t}}-\frac1t-\frac12
       -\sum_{j=1}^{M}\frac{B_{2j}}{(2j)!}t^{2j-1}\right)
 \quad(t>0),
\]
extended by zero for \(t\leq0\). The density tends to zero at the
origin on the stated sufficient half-plane, and its
Fourier coefficients are absolutely summable there. The standard
Hurwitz expansion \cite[Section 25.11]{DLMF} gives the required
normal convergence. We credit \eqref{eq:root-master} to the preceding
report; it is the starting point, not a new claim.

\subsection{A universal power-lattice cutoff lemma}

The following elementary lemma isolates the whole obstruction.

\begin{lemma}\label{lem:power-cutoff}
For \(a,h\) as above and \(u\) in a sufficiently small complex
neighborhood of zero, put
\[
 P_K(u)=\sum_{n=-K}^{K}(a+\ii hn)^{-1-u},\qquad
 P_\infty^{\mathrm{an}}(u)=
       \frac{\tau^{1+u}}{\Gamma(1+u)}\Li_{-u}(e^{-\tau a}).
\]
Then, locally uniformly in \(u\) and \(a\),
\begin{equation}\label{eq:power-cutoff}
 P_K(u)=P_\infty^{\mathrm{an}}(u)
    +\frac{2}{h}\frac{\sin(\pi u/2)}{u}
                   e^{-u\log(hK)}
    +O\!\left(K^{-1-\Ree u}\right).
\end{equation}
The sine quotient is filled by its removable value. Derivatives of
order \(j\) of the error at zero are
\(O(K^{-1}(1+\log K)^j)\).
\end{lemma}

\begin{proof}
For \(\Ree u>0\), ordinary periodization of
\(t^u e^{-at}/\Gamma(1+u)\) gives
\(\sum_n(a+\ii hn)^{-1-u}=P_\infty^{\mathrm{an}}(u)\).
The omitted paired tail has the expansion
\[
 (a+\ii hn)^{-1-u}+(a-\ii hn)^{-1-u}
 =-2\sin(\pi u/2)(hn)^{-1-u}+O(n^{-2-\Ree u}).
\]
Since
\(\sum_{n>K}n^{-1-u}=K^{-u}/u+O(K^{-1-\Ree u})\),
subtracting the tail proves \eqref{eq:power-cutoff} in that half-plane.
To see the uniform continuation through \(u=0\), use the exact tail
expression
\begin{align*}
 &( \ii h)^{-1-u}\zeta(1+u,K+1-\ii a/h)\\
 &\hspace{12mm}
 +(-\ii h)^{-1-u}\zeta(1+u,K+1+\ii a/h).
\end{align*}
Its two residues cancel at \(u=0\). Euler--Maclaurin applied to
these two Hurwitz functions proves the stated expansion on a full
neighborhood. Differentiating that expansion introduces only powers
of \(\log K\). The sector and parameter estimates are uniform on
compact subsets of \(\Ree a>0\).
\end{proof}

At \(u=0\), the lemma gives the familiar distinction
\[
 \lim_{K\to\infty}P_K(0)
 =\frac{\tau}{2}\coth\frac{\tau a}{2}
 =\tau\Li_0(e^{-\tau a})+\frac{\tau}{2}.
\]
The extra \(\tau/2\) is the zeroth member of the entire cutoff
generating function in \eqref{eq:power-cutoff}.

\subsection{The all-order contact polynomial}

For an even nonnegative integer \(N\), define
\begin{equation}\label{eq:root-beta}
 M=N/2,\qquad
 \beta_N=\frac{B_{N+2}}{(N+1)(N+2)},\qquad
 A_N(u)=\prod_{j=1}^{N}\left(1-\frac uj\right).
\end{equation}
An empty product equals one. Define the polynomial
\begin{equation}\label{eq:contact-polynomial}
 \boxed{\displaystyle
 C_{N,d}(X)=
  \frac{\tau\beta_N}{2}\,d!\,[u^d]\!
  \left\{uA_N(u)\frac{\sin(\pi u/2)}{\pi u/2}e^{-Xu}\right\},
  \qquad d\geq1.}
\end{equation}
Its degree is \(d-1\), with nonzero leading coefficient
\((\tau\beta_N/2)d(-1)^{d-1}\).

\begin{theorem}[Complete minimal-contact law]\label{thm:minimal-contact}
For every even \(N\geq0\), every \(d\geq1\), and \(\Ree a>0\),
\begin{equation}\label{eq:minimal-contact}
 \boxed{\displaystyle
 \sum_{n=-K}^{K}
   \left.\partial_s^dR_{N/2}(s,a+\ii hn)\right|_{s=-N}
 =
 {\cal F}_{N/2}^{(d)}(-N;a,h)
 +C_{N,d}\bigl(\log(hK)\bigr)
 +O\!\left(\frac{(1+\log K)^{d-1}}K\right).}
\end{equation}
Here and below the superscript on \({\cal F}\) denotes the order
derivative. The error is uniform for \(a\) in compact subsets of
the right half-plane.

For odd \(N\geq1\), put \(M=(N+1)/2\). Then all the corresponding
derivative sums are absolutely convergent, and
\begin{equation}\label{eq:odd-all-jets}
 \sum_{n\in\mathbb Z}
 \left.\partial_s^dR_M(s,a+\ii hn)\right|_{s=-N}
 ={\cal F}_M^{(d)}(-N;a,h).
\end{equation}
\end{theorem}

\begin{proof}
For even \(N\), the difference of two successive compensations is
exactly
\begin{equation}\label{eq:one-more-subtraction}
 R_{N/2}(-N+u,w)-R_{N/2+1}(-N+u,w)
 =\beta_N uA_N(u)w^{-1-u}.
\end{equation}
Indeed,
\[
 \frac{B_{N+2}}{(N+2)!}(-N+u)_{N+1}
       =\beta_N uA_N(u).
\]
The oversubtracted series is normally convergent in
\(\Ree u>-2\). Its omitted tail after \(d\) derivatives at zero is
\(O(K^{-2}(1+\log K)^{d-1})\), since its leading coefficient
contains the factor \(u\). In the initial half-plane \(\Ree u>0\),
\eqref{eq:root-master} and \eqref{eq:one-more-subtraction} also imply
\[
 {\cal F}_{N/2}(-N+u)-{\cal F}_{N/2+1}(-N+u)
       =\beta_NuA_N(u)P_\infty^{\mathrm{an}}(u).
\]
Both sides continue holomorphically. Apply Lemma~\ref{lem:power-cutoff}
to the truncated version of \eqref{eq:one-more-subtraction}. Its
explicit cutoff contribution is
\[
 \frac{2\beta_N}{h}A_N(u)\sin(\pi u/2)e^{-u\log(hK)}
 =\frac{\tau\beta_N}{2}\,
       uA_N(u)\frac{\sin(\pi u/2)}{\pi u/2}e^{-u\log(hK)}.
\]
Taking the \(d\)th coefficient proves the theorem and the stated
error. For odd \(N\), the minimally chosen \(M\) already satisfies
\(-N>-2M\), so \eqref{eq:odd-all-jets} is an immediate application
of normal convergence in \eqref{eq:root-master}.
\end{proof}

This proof also explains why merely differentiating the first-order
Bernoulli correction is insufficient: the factor
\(A_N(u)\sin(\pi u/2)/(\pi u/2)\) itself contributes higher
coefficients.

\begin{corollary}[Two explicit cutoff conventions]\label{cor:cutoff-coordinate}
The constant term of the cutoff expansion in powers of \(\log K\)
is
\begin{equation}\label{eq:index-finite-part}
 \FP_{\log K}\sum_{n\in\mathbb Z}
     \partial_s^dR_{N/2}(-N,a+\ii hn)
 ={\cal F}_{N/2}^{(d)}(-N;a,h)+C_{N,d}(\log h).
\end{equation}
The constant term when the large variable is instead the physical
radius \(R=hK\) is
\begin{equation}\label{eq:radius-finite-part}
 \FP_{\log R}\sum_{|hn|\leq R}
     \partial_s^dR_{N/2}(-N,a+\ii hn)
 ={\cal F}_{N/2}^{(d)}(-N;a,h)+C_{N,d}(0).
\end{equation}
Only positive powers of the specified logarithm are discarded in
either definition. Their difference is exactly
\(C_{N,d}(\log h)-C_{N,d}(0)\).
\end{corollary}

\begin{proof}
Theorem~\ref{thm:minimal-contact} supplies a polynomial plus a term
that tends to zero. Taking its constant coefficient in the indicated
variable proves both formulas. Replacing \(R/h\) by its integer part
changes the asymptotic by a quantity tending to zero.
\end{proof}

\subsection{Harmonic and zeta coefficients}

Let \(Y_j\) denote the complete exponential Bell polynomials,
defined by
\[
 \exp\left(\sum_{r=1}^{\infty}x_r\frac{u^r}{r!}\right)
   =\sum_{j=0}^{\infty}Y_j(x_1,\ldots,x_j)\frac{u^j}{j!}.
\]
The Euler product for sine and the finite product in
\eqref{eq:root-beta} give
\begin{align}
 x_1&=-X-H_N,\notag\\
 x_r&=-(r-1)!\left(H_N^{(r)}
                +\mathbf1_{2\mid r}\,2^{1-r}\zeta(r)\right),
                  \qquad r\geq2,\label{eq:contact-bell}\\
 C_{N,d}(X)&=\frac{\tau\beta_N}{2}\,d\,
       Y_{d-1}(x_1,\ldots,x_{d-1}).\label{eq:contact-bell-value}
\end{align}
Thus the contact coefficients involve finite harmonic numbers and
even zeta values only; no unspecified Laurent coefficient remains.
Writing \(L=X+H_N\) and
\(V=H_N^{(2)}+\zeta(2)/2\), the first four are
\begin{align}
 C_{N,1}&=\frac{\tau\beta_N}{2},&
 C_{N,2}&=-\tau\beta_N L,\notag\\
 C_{N,3}&=\frac{3\tau\beta_N}{2}(L^2-V),&
 C_{N,4}&=2\tau\beta_N
       \bigl(-L^3+3LV-2H_N^{(3)}\bigr).
 \label{eq:first-four-contacts}
\end{align}
For \(N=0\), in particular,
\begin{equation}\label{eq:gamma-cubic-contact}
 C_{0,1}=\frac{\tau}{24},\qquad
 C_{0,2}(X)=-\frac{\tau X}{12},\qquad
 C_{0,3}(X)=\frac{\tau}{8}
                    \left(X^2-\frac{\pi^2}{12}\right).
\end{equation}
The constant \(-\tau\pi^2/96\) at the third derivative is missed
if the sine factor is dropped.

For explicit evaluation of the analytic term in
\eqref{eq:minimal-contact}, define \(E_j^{(N)}\) by
\begin{equation}\label{eq:root-E}
 \exp\left((\log\tau+\gamma-H_N)u+
   \sum_{r=2}^{\infty}
       \frac{(-1)^{r+1}\zeta(r)-H_N^{(r)}}r u^r\right)
   =\sum_{j=0}^{\infty}E_j^{(N)}\frac{u^j}{j!}.
\end{equation}
The reciprocal-Gamma factorization then gives, for every \(M\),
\begin{equation}\label{eq:root-analytic-jets}
 {\cal F}_M^{(d)}(-N;a,h)
 =(-1)^NN!\tau^{-N}d
   \sum_{j=0}^{d-1}\binom{d-1}{j}E_j^{(N)}
       \partial_s^{d-1-j}{\cal B}_M(-N;a,\tau).
\end{equation}
This coefficient formula is inherited from \cite{RepoVertical};
the new result supplies its exact minimal-cutoff interpretation.
All order derivatives on the right are normally convergent
polylogarithm or Lambert--polylogarithm derivatives.

\subsection{Periodic filters: the mean is the only obstruction}

Let \(c_n\) have period \(p\) and mean
\(\bar c=p^{-1}\sum_{r=0}^{p-1}c_r\). Define
\begin{equation}\label{eq:filtered-analytic}
 {\cal F}_{M,c}(s;a,h)=
    \sum_{r=0}^{p-1}c_r\,
       {\cal F}_M(s;a+\ii hr,ph).
\end{equation}

\begin{theorem}[All-order filter criterion]\label{thm:filter-all}
For even \(N\) and \(d\geq1\),
\begin{equation}\label{eq:weighted-contact}
 \sum_{n=-K}^{K}c_n\partial_s^dR_{N/2}(-N,a+\ii hn)
 =
 {\cal F}_{N/2,c}^{(d)}(-N;a,h)
 +\bar c\,C_{N,d}(\log(hK))+o(1).
\end{equation}
If \(\bar c=0\), the positive and negative halves both converge
ordinarily, locally uniformly in the order parameter near \(-N\);
every derivative equals the corresponding derivative of
\eqref{eq:filtered-analytic}. No finite part is needed.
The condition \(\bar c=0\) is necessary and sufficient for all
contacts to vanish. There are no additional discrete-moment
conditions at higher derivative order.
\end{theorem}

\begin{proof}
Decompose \(c_n=\bar c+(c_n-\bar c)\). The constant part is
Theorem~\ref{thm:minimal-contact}. Partial sums of the mean-zero
part are bounded. Summation by parts, using
\[
 \partial_s^dR_{N/2}(-N,a+\ii hn)
     =O\!\left(\frac{(1+\log|n|)^{d-1}}{|n|}\right)
\]
and its first difference
\(O((1+\log|n|)^{d-1}/n^2)\), proves convergence of each half.
The same argument before differentiation is locally uniform in
\(\Ree(s+N)>-1\). In \(\Ree(s+N)>0\) the sums equal
\eqref{eq:filtered-analytic} by absolute convergence and residue-class
splitting; analytic continuation proves equality near \(s=-N\).
Necessity follows already from
\(C_{N,1}=\tau\beta_N/2\neq0\).
\end{proof}

For the alternating filter this yields the particularly simple
all-order identity
\begin{equation}\label{eq:alternating-all-jets}
 \sum_{n\in\mathbb Z}(-1)^n
      \partial_s^dR_{N/2}(-N,a+\ii hn)
 =
 2{\cal F}_{N/2}^{(d)}(-N;a,2h)
   -{\cal F}_{N/2}^{(d)}(-N;a,h).
\end{equation}
The two half-sums specify the summation convention; absolute
convergence is not asserted. For every two shifts \(a,b\) in the
right half-plane, the relative identity
\begin{equation}\label{eq:relative-lattice}
 \begin{split}
 &\sum_{n\in\mathbb Z}c_n
 \left\{\partial_s^dR_{N/2}(-N,a+\ii hn)
       -\partial_s^dR_{N/2}(-N,b+\ii hn)\right\}\\
 &\hspace{10mm}
 ={\cal F}_{N/2,c}^{(d)}(-N;a,h)
       -{\cal F}_{N/2,c}^{(d)}(-N;b,h)
 \end{split}
\end{equation}
is absolutely convergent for every bounded periodic \(c\).
The leading reciprocal-logarithmic term is independent of the
horizontal shift, so subtraction improves it to
\(O(n^{-2}(1+\log|n|)^{d-1})\).

\subsection{Gamma and all Stieltjes antiderivatives}

At \(N=0\), write \(D_d(w)=\partial_s^dR_0(0,w)\). Direct
differentiation of the fixed subtractions gives
\begin{equation}\label{eq:Dd-explicit}
 D_d(w)=\zeta^{(d)}(0,w)
       +d!\,w\sum_{j=0}^{d}\frac{(-\Log w)^j}{j!}
       -\frac12(-\Log w)^d.
\end{equation}
The first instance is the usual Stirling remainder,
\[
 D_1(w)=\log\Gamma(w)-(w-\tfrac12)\Log w+w-\tfrac12\log(2\pi).
\]
Its vertical Gamma product belongs to the published work of
Arabi Ardehali and Rosengren \cite{ArdehaliRosengren}, as the
preceding report correctly records. Our contribution here is the
explicit all-\(d\) minimal-cutoff and filter law.

For example,
\begin{equation}\label{eq:D2-cutoff}
 \begin{split}
 &\sum_{n=-K}^{K}
 \left\{\zeta''(0,w_n)+w_n\bigl(\Log^2w_n-2\Log w_n+2\bigr)
                       -\tfrac12\Log^2w_n\right\}\\
 &\hspace{9mm}
 ={\cal F}_0''(0;a,h)-\frac{\tau}{12}\log(hK)
            +O(K^{-1}\log K).
 \end{split}
\end{equation}
This identity is simultaneously a formula for compensated
antiderivatives of the first generalized Stieltjes constant:
expansion of
\(\partial_w\zeta(s,w)=-s\zeta(s+1,w)\) gives
\begin{equation}\label{eq:stieltjes-primitive-root}
 \partial_w\zeta^{(d)}(0,w)=(-1)^d d\,\gamma_{d-1}(w),
 \qquad d\geq1.
\end{equation}
Thus
\[
 {\mathfrak A}_m(w)=\frac{(-1)^{m+1}}{m+1}
       \left(\zeta^{(m+1)}(0,w)-\zeta^{(m+1)}(0,1)\right)
\]
is the uniquely anchored primitive
\(\mathfrak A_m'(w)=\gamma_m(w)\),
\(\mathfrak A_m(1)=0\). Theorem~\ref{thm:minimal-contact},
\eqref{eq:Dd-explicit}, and \eqref{eq:root-analytic-jets}
evaluate every corresponding compensated lattice sum using finite
harmonic/zeta coefficients and Lambert--polylogarithm order
derivatives. Formula \eqref{eq:relative-lattice} supplies an
ordinary absolutely convergent version whenever two anchors are
compared.


% END INPUT


% BEGIN sections/03_modular_transport.tex
\section{Exact modular transport of the order derivatives}
\label{sec:modular}

Question 7 of \cite{RepoVertical} asks for an exact
\(\tau\leftrightarrow4\pi^2/\tau\) transport, including the
nonmodular terms. At horizontal shift \(a=1\), we give one
holomorphic identity near every nonpositive even order and another
near the Stieltjes order \(s=1\). Their derivatives are specified
by finite polynomial recurrences and complete residue formulas.

This section uses classical Mellin-transform technology. The eta
and Ramanujan specializations, the zeta functional equation, and
the existence of modular relations for generalized Stieltjes
constants are established literature, not novelty claims.
In particular, the all-index modular relations of Dixit,
Sathyanarayana and Sharan \cite{DixitStieltjes} preclude any claim
that modular Stieltjes identities themselves originate here.
The contribution to the present continuation is the explicit
order-dependent kernel, its residue compiler, and its compatibility
with the fixed-cutoff law of Section~\ref{sec:lattice}.

\subsection{One contour identity before coefficient extraction}

Put
\begin{equation}\label{eq:mod-lambda}
 \Lambda(s,\tau)=\sum_{n=1}^{\infty}
       \frac{n^{s-1}}{e^{n\tau}-1},\qquad
 \tau^*=\frac{4\pi^2}{\tau},\qquad
 \alpha=\frac{2\pi}{\tau},\qquad L=\log\alpha.
\end{equation}
The function \(\Lambda\) is entire in \(s\). It is exactly
\({\cal L}_{1-s}(e^{-\tau},e^{-\tau})\), so subtracting the
finite polylogarithm terms in \eqref{eq:root-bracket} recovers
the vertically sampled family.

For \(M\geq0\), let
\begin{align}
 {\cal Q}_M(s,\tau)
  ={}&\frac{\zeta(2-s)}{\tau}
   +\Gamma(s)\zeta(s)\tau^{-s}-\frac12\zeta(1-s)\notag\\
 &+\sum_{j=1}^{M+1}\frac{B_{2j}}{(2j)!}
          \zeta(2-2j-s)\tau^{2j-1}.
 \label{eq:mod-residue-even}
\end{align}
Near \(s=1\), instead put
\begin{equation}\label{eq:mod-residue-one}
 {\cal Q}_*(s,\tau)=
   \frac{\zeta(2-s)}{\tau}
   +\Gamma(s)\zeta(s)\tau^{-s}-\frac12\zeta(1-s).
\end{equation}
The apparently singular terms are combined before their removable
values are taken.

\begin{theorem}[Local modular transport]\label{thm:modular}
Fix \(\tau>0\). If \(M\geq0\) and \(|s+2M|<1/4\), then
\begin{equation}\label{eq:modular-master}
 \boxed{\displaystyle
 \Lambda(s,\tau)={\cal Q}_M(s,\tau)
  +\frac{\alpha^s}{2\pi\ii}
   \int_{\Ree w=3/2}
       \Gamma(w)\zeta(w)\zeta(w+1-s)
       \frac{\cos(\pi w/2)}{\cos(\pi(s-w)/2)}
       (\tau^*)^{-w}\,\dd w.}
\end{equation}
If \(|s-1|<1/4\), the same identity holds with
\({\cal Q}_M\) replaced by \({\cal Q}_*\).
The contour is oriented upward. Both identities are holomorphic
through their central orders, and the integral admits every fixed
order derivative under its sign.
\end{theorem}

\begin{proof}
For \(c>\max(1,\Ree s)\), Mellin inversion and absolute convergence
give
\begin{equation}\label{eq:lambert-mellin}
 \Lambda(s,\tau)=\frac1{2\pi\ii}
       \int_{\Ree z=c}\Gamma(z)\zeta(z)\zeta(z+1-s)
                                  \tau^{-z}\,\dd z.
\end{equation}
In the stated neighborhoods one may take \(c=3/2\) and move the
contour to \(\Ree z=\Ree s-3/2\).
Near \(s=-2M\), the poles crossed are \(z=1,s,0\) and
\(z=1-2j\), \(1\leq j\leq M+1\). Their residues are exactly
\eqref{eq:mod-residue-even}. The negative even poles of Gamma
are canceled by the trivial zeros of \(\zeta(z)\).
Near \(s=1\), only \(z=1,s,0\) are crossed, giving
\eqref{eq:mod-residue-one}.

The two zeta functional equations \cite[Section 25.4]{DLMF} imply
the meromorphic identity
\begin{equation}\label{eq:mod-functional-factor}
 \begin{split}
 \Gamma(z)\zeta(z)\zeta(z+1-s)
  ={}&(2\pi)^{2z-s}\Gamma(s-z)\zeta(s-z)\zeta(1-z)\\
 &\times
 \frac{\cos(\pi(s-z)/2)}{\cos(\pi z/2)}.
 \end{split}
\end{equation}
Substitute \(w=s-z\) on the new contour; the resulting integral
is the one in \eqref{eq:modular-master}.
Stirling's formula gives exponential decay, while the zeta
factors have at most polynomial growth on the fixed vertical
strips. The trigonometric quotient and all its order derivatives
are bounded there on compact parameter sets: the denominator has
no zero on the contour in either stated neighborhood. These facts
justify the contour motion and differentiated integrals.
Away from the central collisions the residue computation is
ordinary. The formulas below explicitly remove those collisions,
or equivalently one may use the identity theorem and the
holomorphic left-hand side.
\end{proof}

At \(s=-2M\), the trigonometric quotient is the constant
\((-1)^M\). Thus the zeroth derivative has the classical
Ramanujan form
\begin{equation}\label{eq:ramanujan-recovered}
 \Lambda(-2M,\tau)=
 {\cal Q}_M(-2M,\tau)+(-1)^M\alpha^{-2M}
                         \Lambda(-2M,\tau^*).
\end{equation}
For \(M=0\), it becomes the eta transformation
\begin{equation}\label{eq:eta-recovered}
 \Lambda(0,\tau)-\Lambda(0,\tau^*)
   =\frac{\zeta(2)}{\tau}
     +\frac12\log\frac{\tau}{2\pi}-\frac{\tau}{24}.
\end{equation}
For example, at \(M=1\),
\begin{align}
 \Lambda(-2,\tau)
 ={}&\frac{\zeta(4)}{\tau}
   -\frac{\zeta(3)}2\left(1+\frac{\tau^2}{4\pi^2}\right)
   +\frac{\zeta(2)\tau}{12}+\frac{\tau^3}{1440}\notag\\
 &-\frac{\tau^2}{4\pi^2}\Lambda(-2,\tau^*).
 \label{eq:zeta3-recovered}
\end{align}
These are useful normalization checks. The positive
\(\tau^3/1440\) sign in the last polynomial is forced by
\(B_4\zeta(0)>0\).

\subsection{A finite polynomial compiler for every derivative}

Define polynomials \(p_j(X)\) by
\begin{equation}\label{eq:mod-kernel-recurrence}
 p_0(X)=1,\qquad
 p_{j+1}(X)=-(1+X^2)p_j'(X)-Xp_j(X).
\end{equation}
They satisfy
\begin{equation}\label{eq:mod-kernel-generator}
 \frac1{\cos v+X\sin v}
      =\sum_{j=0}^{\infty}p_j(X)\frac{v^j}{j!}.
\end{equation}
Indeed, differentiating a shifted secant proves the recurrence.
The first values are
\[
 p_1=-X,\qquad p_2=1+2X^2,\qquad
 p_3=-5X-6X^3,\qquad p_4=5+28X^2+24X^4.
\]

For \(M,j,\ell\geq0\), set
\begin{equation}\label{eq:Tmod}
 {\cal T}_{j,\ell}^{(M)}(x)=
 \frac{(-1)^\ell}{2\pi\ii}
 \int_{\Ree w=3/2}
  \Gamma(w)\zeta(w)\zeta^{(\ell)}(w+1+2M)
  p_j\!\left(\tan\frac{\pi w}{2}\right)x^{-w}\,\dd w.
\end{equation}
Every integral is absolutely convergent. For \(j=0\),
\[
 {\cal T}_{0,\ell}^{(M)}(x)
       =\left.\partial_s^\ell\Lambda(s,x)\right|_{s=-2M}.
\]

\begin{corollary}[All negative-even modular derivatives]\label{cor:mod-jets}
For every \(r\geq0\),
\begin{equation}\label{eq:mod-jets}
 \boxed{\displaystyle
 \partial_s^r\Lambda(-2M,\tau)
 ={\cal Q}_M^{(r)}(-2M,\tau)
  +(-1)^M\alpha^{-2M}
  \sum_{\substack{j,\ell,k\geq0\\j+\ell+k=r}}
       \frac{r!}{j!\ell!k!}
       \left(\frac{\pi}{2}\right)^j
       L^k\,{\cal T}_{j,\ell}^{(M)}(\tau^*).}
\end{equation}
\end{corollary}

\begin{proof}
Put \(s=-2M+u\). The quotient in
\eqref{eq:modular-master} becomes
\[
 \frac{(-1)^M}
 {\cos(\pi u/2)+\tan(\pi w/2)\sin(\pi u/2)}.
\]
Expand this by \eqref{eq:mod-kernel-generator},
\(\alpha^u=e^{Lu}\), and
\(\zeta(w+1+2M-u)\). Normal convergence allows coefficient
extraction. The multinomial coefficient is the ordinary
Leibniz coefficient.
\end{proof}

The terms \(j=0\) are the familiar transformed Lambert
derivatives. The finitely many \(j\geq1\) terms are the exact
nonmodular correction. Omitting them is already wrong at the
first derivative. This is an identity with explicitly specified
ordinary contour integrals, not an assertion that those integrals
always reduce to a finite list of polylogarithm values.

The corresponding critical-order compiler is equally explicit.
Define
\begin{equation}\label{eq:Umod}
 {\cal U}_{j,\ell}(x)=
 \frac{(-1)^\ell}{2\pi\ii}
 \int_{\Ree w=3/2}
    \Gamma(w)\zeta(w)\zeta^{(\ell)}(w)
    Xp_j(-X)x^{-w}\,\dd w,\qquad
    X=\cot\frac{\pi w}{2}.
\end{equation}
Since the trigonometric quotient at \(s=1+u\) is
\[
 \frac{X}{\cos(\pi u/2)-X\sin(\pi u/2)},
\]
we obtain
\begin{equation}\label{eq:critical-mod-jets}
 \partial_s^r\Lambda(1,\tau)
  ={\cal Q}_*^{(r)}(1,\tau)
   +\alpha\!
    \sum_{j+\ell+k=r}\frac{r!}{j!\ell!k!}
          \left(\frac{\pi}{2}\right)^j L^k
                                  {\cal U}_{j,\ell}(\tau^*).
\end{equation}
Thus every critical-order derivative is transported by finitely
many integrals involving \(\zeta(w)\zeta^{(\ell)}(w)\).

\subsection{The colliding residues and their Stieltjes coefficients}

The finite residue terms in the preceding formulas require care.
At zero, let
\begin{equation}\label{eq:mod-e}
 e_j=\left.\frac{\dd^j}{\dd u^j}
                 \left(e^{Lu}\sec\frac{\pi u}{2}\right)\right|_{u=0}.
\end{equation}
Then \(e_0=1\), \(e_1=L\),
\(e_2=L^2+\pi^2/4\), and
\(e_3=L^3+3\pi^2L/4\). The zeta functional equation gives
\[
 \Gamma(s)\zeta(s)\tau^{-s}-\frac12\zeta(1-s)
 =\frac12\zeta(1-s)
           \left(\alpha^s\sec\frac{\pi s}{2}-1\right).
\]
The factor in parentheses vanishes at zero. Multiplying the
full Laurent series, including its pole, yields
\begin{equation}\label{eq:Qzero-jets}
 \begin{split}
 {\cal Q}_0^{(r)}(0,\tau)
 ={}&(-1)^r\left(\frac{\zeta^{(r)}(2)}{\tau}
                  +\frac{\tau}{12}\zeta^{(r)}(0)\right)\\
 &+\frac12\left\{-\frac{e_{r+1}}{r+1}
        +\sum_{j=0}^{r-1}\binom rj\gamma_j e_{r-j}\right\}.
 \end{split}
\end{equation}
The sum is empty at \(r=0\). In particular, the residue at
derivative order \(r\) needs one extra coefficient \(e_{r+1}\).
Replacing the Laurent germ by its regular part would delete this
term.

For \(M\ge1\) there is no Laurent collision at \(s=-2M\).
The same \(e_j\), with \(L=\log\alpha\), give every residue
derivative without separately differentiating a singular Gamma
factor:
\begin{equation}\label{eq:Qnegative-jets}
 \begin{split}
 {\cal Q}_M^{(r)}(-2M,\tau)
 ={}&(-1)^r\left\{\frac{\zeta^{(r)}(2M+2)}{\tau}
    +\sum_{j=1}^{M+1}\frac{B_{2j}\tau^{2j-1}}{(2j)!}
                    \zeta^{(r)}(2M+2-2j)\right\}\\
 &+\frac12\left\{(-1)^M\alpha^{-2M}
      \sum_{\ell=0}^{r}\binom r\ell(-1)^\ell
                \zeta^{(\ell)}(2M+1)e_{r-\ell}
              -(-1)^r\zeta^{(r)}(2M+1)\right\}.
 \end{split}
\end{equation}
Indeed the functional-equation expression immediately preceding
\eqref{eq:Qzero-jets} is a product of two holomorphic functions
at these orders. Ordinary Leibniz differentiation proves
\eqref{eq:Qnegative-jets}.

At \(s=1\), let
\begin{equation}\label{eq:mod-A}
 a_j=\left.\frac{\dd^j}{\dd u^j}
            \left(\Gamma(1+u)\tau^{-u}\right)\right|_{u=0}.
\end{equation}
These are Bell polynomials generated by
\[
 \Gamma(1+u)\tau^{-u}
 =\exp\left(-(\gamma+\log\tau)u+
          \sum_{m=2}^{\infty}\frac{(-1)^m\zeta(m)}m u^m\right).
\]
The two colliding terms at \(s=1+u\) combine as
\[
 \frac1{\tau}
 \left\{\frac{A(u)-1}{u}
       +\sum_{j=0}^{\infty}\frac{\gamma_j u^j}{j!}
                              \bigl(1+(-1)^j A(u)\bigr)\right\},
 \qquad A(u)=\Gamma(1+u)\tau^{-u}.
\]
Consequently,
\begin{equation}\label{eq:Qone-jets}
 \boxed{\displaystyle
 {\cal Q}_*^{(r)}(1,\tau)=
 \frac1{\tau}\left\{\frac{a_{r+1}}{r+1}+\gamma_r+
     \sum_{j=0}^{r}\binom rj(-1)^j\gamma_j a_{r-j}\right\}
       -\frac{(-1)^r}{2}\zeta^{(r)}(0).}
\end{equation}
This is a finite formula in ordinary Stieltjes constants,
zeta values, and logarithms for every \(r\).
The first two cases are
\begin{align}
 {\cal Q}_*(1,\tau)&=\frac{\gamma-\log\tau}{\tau}+\frac14,
 \label{eq:Qone-first}\\
 {\cal Q}_*'(1,\tau)&=
       \frac{\log^2\tau-\gamma^2+\zeta(2)}{2\tau}
                    +\frac12\zeta'(0).
 \label{eq:Qone-second}
\end{align}
Cancellation of \(\gamma_1\) in the latter identity is exact.

\subsection{A self-dual integral evaluation}

At \(\tau=2\pi\), \(M=0\), and \(r=1\), the Lambert
derivatives on the two sides of \eqref{eq:mod-jets} cancel.
Using \eqref{eq:Qzero-jets} gives the explicit identity
\begin{equation}\label{eq:self-dual-integral}
 \boxed{\displaystyle
 \frac1{2\pi\ii}\int_{\Ree w=3/2}
   \Gamma(w)\zeta(w)\zeta(w+1)
      \tan\frac{\pi w}{2}\,(2\pi)^{-w}\,\dd w
 =-\frac{\zeta'(2)}{\pi^2}-\frac{\pi}{8}
                              +\frac{\log(2\pi)}6.}
\end{equation}
The integral is absolutely convergent; no principal value occurs.
Higher self-dual derivatives give finite linear identities among
the integrals \({\cal T}_{j,\ell}^{(0)}(2\pi)\), with complete
Stieltjes coefficients supplied by \eqref{eq:Qzero-jets}.

Equations \eqref{eq:modular-master}--\eqref{eq:critical-mod-jets},
together with the finite subtractions in \eqref{eq:root-bracket},
answer the modular-transport question at \(a=1\) in the indicated
order neighborhoods. General horizontal shifts, and finite
polylogarithmic reductions of the correction integrals, remain
separate questions.


% END INPUT


% BEGIN sections/04_mixed_tail_theory.tex
% Integration fragment. Requires amsmath, amssymb, amsthm and hyperref.
% All labels and bibliography keys have the prefix mtp:.
\section{A polarized Bernoulli calculus for arbitrary mixed even tails}
\label{mtp:section}

The all-power continuation in \emph{Bernoulli Tails and Spectral Faces}
proves its recurrence for powers of one even-order Hurwitz tail and
explicitly leaves products containing unequal orders among three or
more factors outside its scope \cite[Section 2, after the even-order
extension]{RepoBernoulli}. The earlier mixed-moment theorem treats
two factors \cite[Section 4]{RepoMixed}. We now remove both
restrictions: the tail orders may be arbitrary positive even integers,
and the number of factors is unrestricted. The result is one finite
triangular recurrence, an anchored analytic generating function, and a
fixed finite arithmetic space for all regular centered moments.

The unweighted sum of an arbitrary product of zeta tails, and its
ordered-index interpretation in multiple zeta values, are classical
\cite{Hoffman}. Those sums are used below as explicit anchors.
The additional assertion concerns \emph{every} centered polynomial
weight, including weights requiring polynomial subtraction. The proof
keeps the cutoff finite until its boundary terms have cancelled.
In particular, this is a continuation of a specified repository question;
we do not claim that the classical tail-product anchors are new.

\subsection{Definitions and ordinary realization of the regular moments}

Let $I$ be a nonempty finite set of labels, let $d=|I|$, and assign an
even integer $p_i\ge2$ to each $i\in I$. Equal assigned integers retain
their different labels. Write
\[
 P_S=\sum_{i\in S}p_i,\qquad w_I=P_I-d,\qquad
 T_p(n)=\zeta(p,n+1)=\zeta(p)-H_n^{(p)},
 \qquad F_I(n)=\prod_{i\in I}T_{p_i}(n),
\]
and put $F_\varnothing(n)=1$, $x_n=n+1/2$. The polygamma normalization is
$T_p(n)=\psi^{(p-1)}(n+1)/(p-1)!$ for even $p$.

For $j\ge0$ define
\begin{equation}\label{mtp:eta}
 \eta_{p,0}=\frac1{p-1},\qquad
 \eta_{p,j}=\frac{(p)_{2j-1}}{(2j)!}B_{2j}(1/2)\quad(j\ge1),
 \qquad
 h_{I,j}=\sum_{\substack{j_i\ge0\\\sum_i j_i=j}}
                        \prod_{i\in I}\eta_{p_i,j_i}.
\end{equation}
The centered Hurwitz asymptotics \cite[Section 25.11]{DLMF} give
\begin{equation}\label{mtp:asymptotic}
                F_I(n)\sim\sum_{j\ge0}h_{I,j}x_n^{-w_I-2j}.
\end{equation}
Every $\eta_{p,j}$ has sign $(-1)^j$, including the positive $j=0$
coefficient. Consequently $h_{I,j}$ is nonzero with the same sign.
This is an asymptotic expansion, not a convergent Laurent series at a
fixed finite $n$.

\begin{proposition}[Regular moments as convergent sums]
\label{mtp:ordinary}
The Dirichlet series
\[
                 \mathcal F_I(s)=\sum_{n\ge0}F_I(n)x_n^{-s},
                    \qquad \Re s>1-w_I,
\]
continues meromorphically to the plane. Its poles are precisely
$1-w_I-2j$, with residue $h_{I,j}$, for $j\ge0$. For every integer
$m\ge0$ satisfying $m\equiv d\pmod2$, its regular value
$M_{I,m}=\mathcal F_I(-m)$ is the ordinary absolutely convergent sum
\begin{equation}\label{mtp:ordinary-formula}
 \boxed{M_{I,m}=\sum_{n=0}^{\infty}\left\{
       x_n^mF_I(n)
       -\sum_{0\le j\le(m-w_I)/2}h_{I,j}x_n^{m-w_I-2j}
                                     \right\}.}
\end{equation}
A sum with a negative upper limit is empty.
\end{proposition}

\begin{proof}
Subtract finitely many terms of \eqref{mtp:asymptotic} from the defining
series and add back the corresponding Hurwitz functions
$h_{I,j}\zeta(s+w_I+2j,1/2)$. Increasing the subtraction order produces
normally convergent remainders in successively larger half-planes.
The nonzero coefficients prove that all the listed poles are actual.
Since $w_I\equiv d\pmod2$, none of these poles is $-m$ on the stated
parity. Every subtracted term in \eqref{mtp:ordinary-formula} contributes
$\zeta(-2r,1/2)=0$, $r\ge0$. After the subtraction the summand is
$O(x_n^{-2})$, or decays faster when no subtraction is needed.
This proves both the value and its ordinary convergence.
\end{proof}

For a full asymptotic expansion in the integer cutoff $N$, let
$\operatorname{CT}_N$ mean the coefficient of $N^0(\log N)^0$.
For any integer $k$ put
\begin{equation}\label{mtp:C}
 C_{I,k}=\operatorname{CT}_N\left\{
     \sum_{n=1}^N n^k\bigl(F_I(n-1)-F_I(n)\bigr)+N^kF_I(N)\right\}.
\end{equation}
These coefficients exist by asymptotic subtraction: each tail has an
expansion in integral inverse powers of $N$, and finite sums of such
powers have power--logarithm expansions with a constant coefficient.
The convention for this coefficient is fixed by \eqref{mtp:C}.

The centered Faulhaber polynomial is
\[
 P_m(N)=\sum_{n=0}^{N-1}(n+1/2)^m
       =\frac{B_{m+1}(N+1/2)-B_{m+1}(1/2)}{m+1}.
\]
Finite summation by parts gives
\[
 \sum_{n=0}^{N-1}x_n^mF_I(n)
 =P_m(N)F_I(N)+\sum_{n=1}^NP_m(n)\bigl(F_I(n-1)-F_I(n)\bigr).
\]
All the removed even-power sums in \eqref{mtp:ordinary-formula} are odd
polynomials in $N$, with zero constant coefficient. Therefore
\begin{equation}\label{mtp:moment-transform}
 M_{I,m}=\sum_{k=1}^{m+1}
       \frac{\binom{m+1}{k}}{m+1}B_{m+1-k}(1/2)C_{I,k}.
\end{equation}
Only indices $k\equiv d+1\pmod2$ occur. Define initially as formal series
\begin{equation}\label{mtp:GH}
 G_I(t)=\sum_{\substack{m\ge0\\m\equiv d\ (2)}}M_{I,m}\frac{t^m}{m!},
 \qquad H_I(t)=2\sinh(t/2)G_I(t),\qquad G_\varnothing=H_\varnothing=0.
\end{equation}
The Faulhaber generating function makes \eqref{mtp:moment-transform}
equivalent to
\begin{equation}\label{mtp:HC}
 H_I(t)=\sum_{\substack{k\ge1\\k\equiv d+1\ (2)}}C_{I,k}\frac{t^k}{k!}.
\end{equation}
No interchange with a divergent infinite moment sum is used.

\subsection{The finite polarized identity and the exact boundary parity}

For a polynomial in independent variables $a_i$, write
$\mathscr D=\sum_i h_i\partial_{a_i}$. On a multilinear polynomial this
operator is nilpotent. The scalar Bernoulli identity
\[
 e^z-1=\frac z2(e^z+1)
          -\sum_{r\ge1}\frac{B_{2r}}{(2r)!}z^{2r}(e^z-1)
\]
can therefore be applied to $\prod_{i\in I}a_i$ as a finite polynomial
identity. It is equivalent to the usual generating function
$(z/2)\coth(z/2)=1+\sum_{r\ge1}B_{2r}z^{2r}/(2r)!$.

\begin{lemma}[Polarized Bernoulli difference]
\label{mtp:polarized}
If $\Delta F_J(n)=F_J(n-1)-F_J(n)$, then for every $n\ge1$
\begin{equation}\label{mtp:polarized-formula}
 \boxed{\begin{aligned}
 \Delta F_I(n)={}&\frac12\sum_{i\in I}n^{-p_i}
                         \bigl(F_{I\setminus\{i\}}(n-1)
                              +F_{I\setminus\{i\}}(n)\bigr)\\
 &-\sum_{\substack{S\subsetneq I\\|S|\ge2\text{ even}}}
                   B_{|S|}n^{-P_S}\Delta F_{I\setminus S}(n).
 \end{aligned}}
\end{equation}
Every product on the right contains fewer tail factors.
\end{lemma}

\begin{proof}
Apply the displayed scalar identity with
$a_i=T_{p_i}(n)$ and $h_i=n^{-p_i}$. Since
$T_{p_i}(n-1)=T_{p_i}(n)+n^{-p_i}$, the exponential of $\mathscr D$
implements the simultaneous shift. The identity
\[
 \frac{\mathscr D^r}{r!}\prod_{i\in I}a_i
  =\sum_{\substack{S\subseteq I\\|S|=r}}
                  \left(\prod_{i\in S}h_i\right)\prod_{i\notin S}a_i
\]
accounts for every coefficient. The term $S=I$ would be the difference
of two constants and is zero. Labels, rather than distinct numerical
orders, supply the multiplicities when some $p_i$ coincide.
\end{proof}

Let $\mathcal A_p(u)$ be the formal expansion of $T_p(N)$ at $u=N^{-1}$:
\begin{equation}\label{mtp:unshifted}
 \mathcal A_p(u)=\frac{u^{p-1}}{p-1}-\frac{u^p}{2}
       +\sum_{r\ge1}\frac{B_{2r}(p)_{2r-1}}{(2r)!}u^{p+2r-1}.
\end{equation}
Put $b_I(k)=[u^k]\prod_{i\in I}\mathcal A_{p_i}(u)$, with
$b_\varnothing(k)=\delta_{k,0}$, and define the odd formal series
$E_i(u)=\mathcal A_{p_i}(u)+u^{p_i}/2$.

\begin{lemma}[Boundary cancellation for unequal increments]
\label{mtp:boundary}
For $k\equiv d+1\pmod2$,
\begin{equation}\label{mtp:boundary-coefficient}
 b_I(k)+\frac12\sum_{i\in I}b_{I\setminus\{i\}}(k-p_i)
 +\sum_{\substack{S\subseteq I\\|S|\ge2\text{ even}}}
                    B_{|S|}b_{I\setminus S}(k-P_S)=0.
\end{equation}
More precisely, the entire generating expression on the left is
\begin{equation}\label{mtp:boundary-polynomial}
 \boxed{\sum_{\substack{S\subseteq I\\|S|\text{ even}}}
             B_{|S|}(1/2)u^{P_S}\prod_{i\notin S}E_i(u).}
\end{equation}
It has parity $d$.
\end{lemma}

\begin{proof}
Let $B_r^+=B_r(1)$, so $B_1^+=1/2$, $B_r^+=B_r$ for $r\ne1$,
and odd indices above one vanish. Before taking a coefficient the
left side of \eqref{mtp:boundary-coefficient} is
\[
                 \sum_{S\subseteq I}B_{|S|}^+u^{P_S}
                              \prod_{i\notin S}\mathcal A_{p_i}(u).
\]
Substitute $\mathcal A_{p_i}=E_i-u^{p_i}/2$. For a fixed subset $R$,
the coefficient of $u^{P_R}\prod_{i\notin R}E_i$ is
\[
 \sum_{S\subseteq R}B_{|S|}(1)(-1/2)^{|R|-|S|}
      =B_{|R|}(1/2),
\]
by the Bernoulli translation identity \cite[24.4.12]{DLMF}.
Odd $|R|$ vanish. Each $u^{P_R}$ is even and the remaining product
has parity $d-|R|\equiv d$. This proves the identity and cancellation.
\end{proof}

This is the point where equal tail orders are unnecessary. Each tail
has its own even increment $u^{p_i}$; Bernoulli translation cancels the
entire mixed boundary polynomial at once. If $d$ is even, the full
subset in \eqref{mtp:boundary-coefficient} has a possible term only at
$k=P_I$, an even index outside the required odd parity. Thus its
inclusion is harmless, although the corresponding difference in
\eqref{mtp:polarized-formula} is omitted.

\subsection{Every anchor is a specified convergent multiple zeta value}

Our convention is
\[
 \zeta(s_1,\ldots,s_\ell)
   =\sum_{n_1>\cdots>n_\ell\ge1}n_1^{-s_1}\cdots n_\ell^{-s_\ell}.
\]
An ordered set partition of $J$ is a sequence
$\pi=(B_1,\ldots,B_\ell)$ of disjoint nonempty blocks whose union is
$J$. For $r\ge1$ and nonempty $J$ define
\begin{align}
 \mathcal T_{J,r}
   &=\sum_{\pi\models J}\zeta(P_{B_1},\ldots,P_{B_\ell},r),
       &\mathcal T_{\varnothing,r}&=\zeta(r)\quad(r\ge2),
       \label{mtp:Tanchor}\\
 \mathcal S_{J,r}
   &=2\mathcal T_{J,r}
       +\sum_{\varnothing\ne U\subseteq J}\mathcal T_{J\setminus U,r+P_U}.
       \label{mtp:Sanchor}
\end{align}
The only use of $J=\varnothing$ below has $r$ even and at least two.
For nonempty $J$ and $h\le0$ put
\begin{equation}\label{mtp:Zanchor}
 \mathcal Z_{J,h}=\sum_{\pi\models J}
       \zeta(P_{B_1},\ldots,P_{B_{\ell-1}},P_{B_\ell}-h),
 \qquad Z_J=\prod_{i\in J}\zeta(p_i),\quad Z_\varnothing=1.
\end{equation}
All these MZVs are convergent and have positive integer entries with
first entry at least two. In particular no isolated $\zeta(1)$ is
introduced.

Index ordering, or direct partition according to the indices equal to
$n$, proves the ordinary identities
\begin{align}
 \mathcal T_{J,r}&=\sum_{n\ge1}n^{-r}F_J(n),\notag\\
 \mathcal S_{J,r}&=\sum_{n\ge1}n^{-r}\bigl(F_J(n-1)+F_J(n)\bigr),
       \label{mtp:anchor-series}\\
 \mathcal Z_{J,h}&=\sum_{n\ge1}n^h\Delta F_J(n)=C_{J,h}\quad(h\le0).
       \notag
\end{align}
The last equality uses $N^hF_J(N)\to0$. These are instances of the
classical tail-index method \cite{Hoffman}; their explicit form
makes every constant in the new recurrence independently computable.
Their respective weights are $P_J+r$, $P_J+r$, and $P_J-h$.

\subsection{The all-multiorder recurrence}

For $m\ge0$ of the regular parity for a nonempty $J$, write
\[
 \mathcal V_J(m)=\sum_{\substack{0\le r\le m\\r\equiv m\ (2)}}
                   \binom mr2^{r-m}M_{J,r}
          =m![t^m]\bigl(\cosh(t/2)G_J(t)\bigr),
 \qquad \mathcal V_\varnothing(m)=0.
\]
For a label $i\in I$, with $J=I\setminus\{i\}$, define
\begin{equation}\label{mtp:W}
 \mathcal W_i(k)=
 \begin{cases}
  \frac12\mathcal S_{J,p_i-k},&k<p_i,\\
  \mathcal V_J(k-p_i)-\frac12\delta_{k,p_i}Z_J,&k\ge p_i.
 \end{cases}
\end{equation}

\begin{theorem}[Triangular mixed-tail recurrence]
\label{mtp:recurrence}
For every $k\ge1$ with $k\equiv d+1\pmod2$,
\begin{equation}\label{mtp:recurrence-formula}
 \boxed{C_{I,k}=\sum_{i\in I}\mathcal W_i(k)
    -\sum_{\substack{S\subsetneq I\\|S|\ge2\text{ even}}}
                     B_{|S|}C_{I\setminus S,k-P_S}.}
\end{equation}
An index $k-P_S\le0$ means the explicit anchor
$\mathcal Z_{I\setminus S,k-P_S}$. Together with
\eqref{mtp:moment-transform}, this is a terminating exact algorithm
for every regular moment, because every right-side tail product has
strictly smaller cardinality.
\end{theorem}

\begin{proof}
Multiply \eqref{mtp:polarized-formula} by $n^k$, sum through a finite
$N$, and then take constant coefficients. When $m=k-p_i<0$, the
symmetric sum is absolutely convergent and equals
$\mathcal S_{J,-m}$. When $m\ge0$, changing the index in one of the
two finite sums gives
\begin{align*}
 \sum_{n=1}^N n^m\bigl(F_J(n-1)+F_J(n)\bigr)
 ={}&\sum_{n=0}^{N-1}\bigl((n+1)^m+n^m\bigr)F_J(n)\\
    &-\delta_{m,0}Z_J+N^mF_J(N).
\end{align*}
The centered polynomial $(x+1/2)^m+(x-1/2)^m$ has the regular parity
for $J$, so the constant coefficient of the first term is
$2\mathcal V_J(m)$. This remains correct for $J=\varnothing$, where
that coefficient is zero by finite power summation.
The difference terms yield $C_{I\setminus S,k-P_S}$ minus their
boundary coefficients. The complete leftover boundary is precisely
\eqref{mtp:boundary-coefficient} and vanishes. What remains is
\eqref{mtp:recurrence-formula}. In particular the exceptional endpoint
term $-Z_J/2$ at $k=p_i$ has not been discarded.
\end{proof}

Let $I^q$ denote the anchored ordinary $q$-fold primitive
\[
 (I^qf)(t)=\frac1{(q-1)!}\int_0^t(t-u)^{q-1}f(u)\,du.
\]
Here the symbol $I^q$ is an integration operator, distinguished from
the label set $I$ by its superscript. Define the finite polynomial
\begin{align}
 R_I(t)={}&\frac12\sum_{i\in I}
       \sum_{\substack{1\le k<p_i\\k\equiv d+1\ (2)}}
          \mathcal S_{I\setminus\{i\},p_i-k}\frac{t^k}{k!}
       -\frac{\mathbf1_{d\ \mathrm{odd}}}{2}
          \sum_{i\in I}Z_{I\setminus\{i\}}\frac{t^{p_i}}{p_i!}
       \notag\\
 &-\sum_{\substack{S\subsetneq I\\|S|\ge2\text{ even}}}B_{|S|}
       \sum_{\substack{1\le k\le P_S\\k\equiv d+1\ (2)}}
          \mathcal Z_{I\setminus S,k-P_S}\frac{t^k}{k!}.
       \label{mtp:R}
\end{align}

\begin{theorem}[Anchored generators, finite arithmetic span, and polylogarithms]
\label{mtp:generator}
The series in \eqref{mtp:GH} converge for $|t|<2\pi$, and are given by
\begin{equation}\label{mtp:generator-formula}
 \boxed{\begin{aligned}
 H_I(t)={}&R_I(t)
       +\sum_{i\in I}I^{p_i}\bigl(\cosh(t/2)G_{I\setminus\{i\}}(t)\bigr)\\
       &-\sum_{\substack{S\subsetneq I\\|S|\ge2\text{ even}}}
                     B_{|S|}I^{P_S}H_{I\setminus S}(t),
       \qquad G_I(t)=\frac{H_I(t)}{2\sinh(t/2)}.
 \end{aligned}}
\end{equation}
For fixed $I$, every $M_{I,m}$ lies in the rational span of $1$ and
convergent positive-index MZVs of weight at most
\begin{equation}\label{mtp:weight-bound}
                         P_I-1-\mathbf1_{d\ \mathrm{odd}}.
\end{equation}
The spanning list is finite and independent of $m$.
Moreover $H_I$ has a finite representation by polynomials in $t$,
MZV constants, and one-variable multiple polylogarithms at $e^{-t}$
with functional depth at most $d-1$. This is a representation bound,
not a claim about minimal depth or arithmetic independence.
\end{theorem}

\begin{proof}
Summing \eqref{mtp:recurrence-formula} with $t^k/k!$ gives
\eqref{mtp:generator-formula} formally. The ranges $k<p_i$, the
endpoint $k=p_i$, and the shifted indices $k-P_S\le0$ give exactly
\eqref{mtp:R}. The remaining positive shifted indices give the displayed
anchored primitives.

Induction on $d$ proves analyticity. All numerators are holomorphic in
$|t|<2\pi$, have parity $d+1$, and vanish to order at least one at
even $d$, or at least two at odd $d$. The denominator has a simple
zero at zero and no other zero in this disk. Its singularity is
therefore removable, and the formal moment series converges there.

Every non-endpoint anchor in $R_I$ has weight $P_I-k$, and its least
possible positive $k$ is one at even $d$ and two at odd $d$. An
endpoint product has weight $P_I-p_i\le P_I-2$ and only occurs at
odd $d$; expand it into MZVs by ordered set partitions. All other
arithmetic constants already belong to generators on proper subsets.
The remaining Taylor coefficients, integrations, and division by
$2\sinh(t/2)$ have rational coefficients. This proves the sharper
parity-sensitive bound \eqref{mtp:weight-bound}, as well as finiteness.
The list can be redundant.

For the functional statement, start at real $t>0$, put $z=e^{-t}$,
and use
\[
 \cosh(t/2)G_J(t)=\left(\frac12+\frac z{1-z}\right)H_J(t).
\]
For $L_{\boldsymbol a}(t)=\operatorname{Li}_{a_1,\ldots,a_\ell}(e^{-t})$
use the convention with $e^{-tn_1}$ on the largest summation index.
Maintain the induction invariant $a_1\ge2$ for nonempty words.
Ordinary integration raises that first index. An explicit unanchored
$q$-fold primitive of $t^m zL_{\boldsymbol a}(t)/(1-z)$ is
\begin{equation}\label{mtp:polylog-primitive}
 (-1)^q\sum_{j=0}^m\binom{q+j-1}{j}\frac{m!}{(m-j)!}
                         t^{m-j}L_{q+j,\boldsymbol a}(t).
\end{equation}
It follows by differentiation from
$L_{1,\boldsymbol a}'=-zL_{\boldsymbol a}/(1-z)$ and
$L_{r+1,\boldsymbol a}'=-L_{r,\boldsymbol a}$. Here $q=p_i\ge2$,
so the new first index is at least two, and depth increases by at
most one.

Subtract the Taylor polynomial of degree $q-1$ at zero from the
\emph{assembled} primitive. Endpoint values with first index at least
two are convergent MZVs. The only possible logarithm comes from a
single first index one. For admissible $\boldsymbol a$,
\begin{align*}
 \lim_{t\to0^+}\bigl(L_{1,\boldsymbol a}(t)
                  -\zeta(\boldsymbol a)L_1(t)\bigr)
 ={}&-\sum_{i=1}^{\ell}\zeta(a_1,\ldots,a_i,1,a_{i+1},\ldots,a_\ell)\\
    &-\sum_{i=1}^{\ell}\zeta(a_1,\ldots,a_i+1,\ldots,a_\ell).
\end{align*}
This follows by subtracting $\zeta(\boldsymbol a)$ inside the outer
series and ordering its missing tail; the difference is absolutely
summable. For the empty word the difference is zero. The coefficient
of $L_1(t)$ in the full endpoint expression is a multiple of $H_J(0)$,
which is zero. Every integration constant is consequently an explicit
convergent MZV. The singleton generator is a polynomial and has depth
zero; induction gives the bound $d-1$. The combined formula is then
continued by its proved holomorphic generator to the stated disk.
\end{proof}

\paragraph{What has been resolved.}
The theorem handles every product of unequal positive even-order tails
and every regular centered moment. It also sharpens the unrestricted
weight bound by one at odd numbers of factors. It neither changes the
summation convention nor assumes a conjecture about MZV independence.
It does not prove the separate all-unequal-order \emph{bilinear optimal
prefix} conjecture: that problem asks for the nonvanishing of different
specified coefficient matrices.

% END INPUT


% BEGIN sections/05_mixed_tail_examples.tex
\section{A mixed cubic, a mixed quartic, and two short convergent identities}
\label{mtp:examples}

In this section a subscript lists the tail orders: $H_{224}$ means
$H_{(2,2,4)}$, whereas $H_4$ is the generator of one fourth-order tail.
Put
\begin{equation}\label{mtp:E5}
 E_5=\frac{15}{2}\zeta(5)-3\zeta(2)\zeta(3)
     =\zeta(2,3)+\zeta(4,1)+\zeta(5).
\end{equation}
The second expression follows from the classical Euler reductions for
odd-weight double zeta values \cite[Equations (4), (9)]{Hoffman}.
The lower generators supplied by Theorem~\ref{mtp:generator} are
\begin{align}
 H_2&=-\frac{t^2}{4},
 &H_4&=\frac{\zeta(2)t^2}{2}-\frac{t^4}{48},\label{mtp:single-generators}\\
 H_{22}&=3\zeta(3)t+2I^2\bigl(\cosh(t/2)G_2\bigr),\notag\\
 H_{24}&=E_5t+\frac{\zeta(3)t^3}{4}
       +I^2\bigl(\cosh(t/2)G_4\bigr)
       +I^4\bigl(\cosh(t/2)G_2\bigr),\label{mtp:lower-generators}\\
 H_{222}&=3I^2\bigl(\cosh(t/2)G_{22}\bigr)
      -\frac{17}{8}\zeta(4)t^2-\frac{\zeta(2)t^4}{48}
      +\frac{t^6}{2880}.\notag
\end{align}
As throughout, $G_J=H_J/(2\sinh(t/2))$. These lower formulas are
inherited special cases, included to specify the subsequent mixed
formulas without unevaluated initial data.

\subsection{All moments of \texorpdfstring{$T_2^2T_4$}{T2 squared times T4}}

\begin{theorem}[A six-coordinate mixed cubic]
\label{mtp:triple}
The complete mixed-cubic generator is
\begin{equation}\label{mtp:H224}
 \boxed{\begin{aligned}
 H_{224}(t)={}&2I^2\bigl(\cosh(t/2)G_{24}\bigr)
                 +I^4\bigl(\cosh(t/2)G_{22}\bigr)\\
       &-\frac5{16}\zeta(6)t^2-\frac7{96}\zeta(4)t^4
                    -\frac{\zeta(2)t^6}{1440}+\frac{t^8}{161280}.
 \end{aligned}}
\end{equation}
Consequently every odd moment $M_{224,2r+1}$ belongs to the rational
span of
\begin{equation}\label{mtp:triple-coordinates}
                         1,\ \zeta(2),\ \zeta(3),\ \zeta(4),\ E_5,\ \zeta(6).
\end{equation}
The first six odd moments have full rank on these six formal labels.
Thus every later odd moment is a rational linear combination of
$M_{224,1},M_{224,3},\ldots,M_{224,11}$.
\end{theorem}

\begin{proof}
Apply \eqref{mtp:generator-formula} to the three labeled orders
$(2,2,4)$. Its nonpolynomial terms are
\[
 2I^2(\cosh G_{24})+I^4(\cosh G_{22})
                      -\frac16I^4H_4-\frac13I^6H_2.
\]
Here and only in this intermediate line, $\cosh$ abbreviates
$\cosh(t/2)$. The anchor polynomial from \eqref{mtp:R} is
\[
          -\frac5{16}\zeta(6)t^2-\frac7{96}\zeta(4)t^4
                                      -\frac{\zeta(2)t^6}{2160}.
\]
For example, its quadratic coefficient simplifies using
\[
 \mathcal S_{\{2,2\},2}
  =4\zeta(2,2,2)+2\zeta(4,2)+2\zeta(2,4)+\zeta(6),
\]
$\zeta(2)\zeta(4)=\zeta(2,4)+\zeta(4,2)+\zeta(6)$, and
$\zeta(2,2,2)=3\zeta(6)/16$. The quartic coefficient uses
$\zeta(2)^2=5\zeta(4)/2$. Inserting the two singleton generators
\eqref{mtp:single-generators} gives exactly \eqref{mtp:H224}.
The lower formulas \eqref{mtp:lower-generators} then prove the asserted
six-coordinate span for every Taylor coefficient.

In the displayed coordinate order, the first six odd moments have
coefficient matrix
\begin{equation}\label{mtp:triple-matrix}
\begingroup\setlength{\arraycolsep}{3pt}\renewcommand{\arraystretch}{1.3}
\begin{pmatrix}
0&0&0&0&1&-5/16\\
0&1/12&1&-7/16&-1/6&5/64\\
-19/360&-29/360&-7/12&35/96&47/720&-35/768\\
761/20160&2257/20160&61/120&-343/768&-257/5040&155/3072\\
-967/25200&-71059/302400&-2419/3360&217/256&9521/134400&-381/4096\\
298871/3193344&3264013/4561920&128981/80640&-9779/4096&-40063/253440&12775/49152
\end{pmatrix}.
\endgroup
\end{equation}
Its exact rational determinant is
\begin{equation}\label{mtp:triple-determinant}
                         -\frac{309971}{13243436236800}\ne0.
\end{equation}
The entries follow by finite Taylor extraction from
\eqref{mtp:H224}; ordinary rational elimination proves
\eqref{mtp:triple-determinant}. Both operations are independently
replayed in the supplied verifier. Inverting this finite matrix proves
the reconstruction assertion. It asserts rank on specified formal
coordinates, not arithmetic independence of the evaluated constants.
\end{proof}

For the first and the first subtracted moments, respectively, this gives
\begin{align}
 \sum_{n\ge0}x_nT_2(n)^2T_4(n)
  &=E_5-\frac5{16}\zeta(6),\label{mtp:triple-m1}\\
 \sum_{n\ge0}\left\{x_n^5T_2(n)^2T_4(n)-\frac13\right\}
  &=-\frac{19}{360}-\frac{29}{360}\zeta(2)-\frac7{12}\zeta(3)
       +\frac{35}{96}\zeta(4)\notag\\
  &\qquad+\frac{47}{720}E_5-\frac{35}{768}\zeta(6).
       \label{mtp:triple-m5}
\end{align}
The subtraction in the second line is fixed by
\[
 T_2(n)^2T_4(n)=\frac1{3x_n^5}-\frac2{9x_n^7}
                    +\frac{211}{1080x_n^9}-\frac{337}{1260x_n^{11}}
                    +O(x_n^{-13}).
\]
In polygamma notation its product is
$\psi'(n+1)^2\psi'''(n+1)/6$.

\begin{corollary}[A short cubic identity]
\label{mtp:short-cubic}
The following series is ordinary and absolutely convergent:
\begin{equation}\label{mtp:short-cubic-formula}
 \boxed{\sum_{n=0}^{\infty}T_2(n)T_4(n)
                     \bigl(1-x_nT_2(n)\bigr)=\frac5{16}\zeta(6).}
\end{equation}
Its summand is $O(x_n^{-6})$.
\end{corollary}

\begin{proof}
The classical unweighted two-tail evaluation gives
$\sum_{n\ge0}T_2(n)T_4(n)=E_5$; it follows directly by sorting the
two indices and is also the linear coefficient in
\eqref{mtp:lower-generators}. Subtract \eqref{mtp:triple-m1}.
Since $T_2(n)=x_n^{-1}-(12x_n^3)^{-1}+O(x_n^{-5})$, the stated
decay follows directly.

There is also a finite telescoping proof, independent of the moment
recurrence. Define
\[
 J(N)=-\frac{N^2}{2}T_2(N)^2T_4(N)+\frac16T_2(N)^3
            -\frac12T_2(N)T_4(N)-\frac16T_6(N).
\]
Direct expansion using $T_p(N-1)=T_p(N)+N^{-p}$ gives the exact identity
\begin{equation}\label{mtp:cubic-telescoper}
 J(N)-J(N-1)=T_2(N-1)T_4(N-1)
                    \bigl(1-(N-1/2)T_2(N-1)\bigr),\qquad N\ge1.
\end{equation}
Every term in $J(N)$ tends to zero. At the other endpoint,
\[
 J(0)=\frac16\zeta(2)^3-\frac12\zeta(2)\zeta(4)-\frac16\zeta(6)
             =-\frac5{16}\zeta(6).
\]
Thus finite summation of \eqref{mtp:cubic-telescoper} gives
\eqref{mtp:short-cubic-formula}. In particular, neither a conjectural
arithmetic relation nor a regularized limiting argument enters this
short evaluation.
\end{proof}

A purely rational identity is also available. Put
\begin{align*}
 P(x)=x\bigl(&2526903+80590275x^2+474388272x^4\\
             &+753887520x^6+353617920x^8+42197760x^{10}\bigr),\\
 Q(x)={}&14065920x^6+108495360x^4+180958272x^2+\frac{338793712}{7}.
\end{align*}
Then
\begin{equation}\label{mtp:rational-sum}
 \sum_{n=0}^{\infty}\bigl\{P(x_n)T_2(n)^2T_4(n)-Q(x_n)\bigr\}
                                      =-6199420.
\end{equation}
The six integer coefficients of $P$ multiply the matrix
\eqref{mtp:triple-matrix} to $(-6199420,0,0,0,0,0)$.
The displayed centered asymptotic coefficients give exactly $Q$;
the residual summand is $O(x_n^{-2})$. Thus this identity follows from
finite rational arithmetic and an ordinary convergent sum.

\subsection{All moments of \texorpdfstring{$T_2^3T_4$}{T2 cubed times T4}}

Define the convergent, weight-nine anchor
\begin{equation}\label{mtp:A9}
 A_9=\sum_{n=0}^{\infty}T_2(n)^3T_4(n)
   =\sum_{\pi\models\{1,2,3,4\}}
        \zeta(P_{B_1},\ldots,P_{B_{\ell-1}},P_{B_\ell}-1),
 \qquad(p_1,p_2,p_3,p_4)=(2,2,2,4).
\end{equation}
The ordered-partition sum is finite; its grouped fifteen-term expansion
is included in the exact verification record. The equality is the
classical minimum-index tail identity, with the $n=0$ term included.
Set
\begin{equation}\label{mtp:B7}
 B_7=\frac{59}{24}\zeta(3)\zeta(4)
              +\frac{23}{12}\zeta(2)\zeta(5)-\frac{203}{32}\zeta(7).
\end{equation}

\begin{theorem}[An eight-coordinate mixed quartic]
\label{mtp:quartic}
The complete even-moment generator is
\begin{equation}\label{mtp:H2224}
 \boxed{\begin{aligned}
 H_{2224}(t)={}&A_9t+B_7t^3-\frac{E_5t^5}{180}\\
 &+3I^2\bigl(\cosh(t/2)G_{224}\bigr)
       +I^4\bigl(\cosh(t/2)G_{222}\bigr)\\
 &-\frac12I^4H_{24}-\frac12I^6H_{22}.
 \end{aligned}}
\end{equation}
Every even moment is therefore in the rational span of the eight
specified coordinates
\begin{equation}\label{mtp:quartic-coordinates}
                1,\ \zeta(2),\ \zeta(3),\ \zeta(4),\ \zeta(6),\ E_5,\ B_7,\ A_9.
\end{equation}
The first eight even moments, of orders $0,2,\ldots,14$, have full
formal rank eight. Their coefficient determinant in this order is
\begin{equation}\label{mtp:quartic-determinant}
                      -\frac{806331520717}{2249069145523200000}\ne0.
\end{equation}
Every later even moment is an explicit rational linear combination
of these eight moments.
\end{theorem}

\begin{proof}
There are three singleton deletions of order two and one of order four.
Among the even proper subsets there are three pairs of type $(2,2)$
and three of type $(2,4)$; multiplication by $B_2=1/6$ gives the two
coefficients $1/2$ in \eqref{mtp:H2224}. The anchor polynomial has
degrees one, three, and five. Its linear coefficient is $A_9$.
Its cubic coefficient is the explicit finite combination
\begin{align}
 B_7={}&\zeta(2,2,2,1)+\frac12\zeta(4,2,1)+\frac12\zeta(2,4,1)
          +\frac12\zeta(2,2,3)\notag\\
       &+\frac16\zeta(6,1)+\frac16\zeta(4,3)-\frac1{12}\zeta(7).
       \label{mtp:B7-MZV}
\end{align}
For completeness, let
\[
 Q_7=\sum_{n\ge0}T_2(n)^4
     =51\zeta(3)\zeta(4)+10\zeta(2)\zeta(5)-\frac{301}{4}\zeta(7).
\]
This is Hoffman's classical quartic evaluation, after adding the
$n=0$ term to his convention \cite[following Corollary 4.4]{Hoffman}.
Sorting its four tail indices shows
\[
 24B_7=Q_7-4\zeta(2,5)-2\zeta(4,3)-3\zeta(7).
\]
Euler's double-zeta evaluations
\[
 \zeta(2,5)=10\zeta(7)-4\zeta(2)\zeta(5)-2\zeta(3)\zeta(4),
 \quad \zeta(4,3)=17\zeta(7)-10\zeta(2)\zeta(5)
\]
give \eqref{mtp:B7}. Finally, the quintic anchor coefficient is
$-(2\zeta(2,3)+\zeta(5))/240=-E_5/180$.
This proves the generator identity and the eight-coordinate span.

The first eight coefficient rows are obtained by Taylor extraction
from the displayed finite system
\eqref{mtp:single-generators}, \eqref{mtp:lower-generators},
\eqref{mtp:H224}, and \eqref{mtp:H2224}. The exact record gives all
$64$ rational entries; the standalone verifier constructs them and
computes \eqref{mtp:quartic-determinant} by rational elimination.
This is a finite nonvanishing certificate. Since the generator already
proves the span for every order, this one exact determinant proves
reconstruction of every later moment from the stated prefix.
No independence of the actual period values is required.
\end{proof}

The first three even moments, in compact form, are
\begin{align}
 M_{2224,0}&=A_9,\notag\\
 M_{2224,2}&=E_5-\frac5{16}\zeta(6)+2B_7-\frac{A_9}{12},
       \label{mtp:quartic-first}\\
 M_{2224,4}&=\frac{\zeta(2)}{20}+\frac32\zeta(3)-\frac{11}{16}\zeta(4)
             +\frac1{16}\zeta(6)-\frac{23}{60}E_5-B_7+\frac7{240}A_9.
       \notag
\end{align}
All three are ordinary convergent sums without subtraction, since
$T_2^3T_4=O(x_n^{-6})$.

\begin{corollary}[A short mixed quartic--cubic identity]
\label{mtp:short-quartic}
One has the ordinary convergent evaluation
\begin{equation}\label{mtp:short-quartic-formula}
 \boxed{\begin{aligned}
 &\sum_{n=0}^{\infty}T_2(n)^2T_4(n)
       \left\{\left(x_n^2+\frac1{12}\right)T_2(n)-x_n\right\}\\
 &\qquad=\frac{59}{12}\zeta(3)\zeta(4)
                +\frac{23}{6}\zeta(2)\zeta(5)-\frac{203}{16}\zeta(7).
 \end{aligned}}
\end{equation}
The summand is $O(x_n^{-8})$.
\end{corollary}

\begin{proof}
The series on the left is
$M_{2224,2}+M_{2224,0}/12-M_{224,1}$. Equations
\eqref{mtp:triple-m1} and \eqref{mtp:quartic-first} reduce it to $2B_7$,
which is the displayed right side. The centered expansion gives
\[
 \left(x^2+\frac1{12}\right)T_2(n)-x
       =\frac1{45x^3}+O(x^{-5}),\qquad x=x_n.
\]
Multiplication by $T_2^2T_4=O(x^{-5})$ proves the asserted decay.
\end{proof}

\subsection{Independent verification and further questions}

The exact verifier uses two different constructions. The first
implements Theorem~\ref{mtp:recurrence}. The second starts from
ordinary finite summation by parts. If
\[
 \mathcal T_{I,q}^{\mathrm{cut}}
  =\operatorname{CT}_N\sum_{n=1}^Nn^qF_I(n),\qquad
       \sum_{n=1}^Nn^q=\sum_a f_{q,a}N^a,
\]
then for $q\ge0$ its independent recurrence is
\begin{equation}\label{mtp:independent-cutoff}
 \mathcal T_{I,q}^{\mathrm{cut}}
   =\sum_a f_{q,a}b_I(a)
    +\sum_{\varnothing\ne S\subseteq I}
       \left\{\sum_a f_{q,a}
                    \mathcal T_{I\setminus S,a-P_S}^{\mathrm{cut}}
                    -\mathcal T_{I\setminus S,q-P_S}^{\mathrm{cut}}\right\}.
\end{equation}
For $q<0$ and nonempty $I$, this is the convergent anchor
$\mathcal T_{I,-q}$. The empty-product cases are zero for $q\ge0$,
Euler's constant for $q=-1$, and $\zeta(-q)$ for $q\le-2$.
The finite identity uses $F_q(n-1)=F_q(n)-n^q$, so its cutoff constants
have exactly the required convention.

For all multisets of one through four orders in $\{2,4,6\}$, the two
routes agree in $306$ regular moments through order $17$, coefficient
by coefficient on unreduced MZV words. Thus no unproved MZV relation
is used to make the comparison succeed. Every Euler-constant label
cancels. There are another $825$ exact boundary-parity checks,
$29$ checks of the simplified mixed generators, the six- and
eight-coordinate determinant certificates, and the rational row in
\eqref{mtp:rational-sum}. The finite cubic telescoper is checked
independently as an exact identity in four Laurent-polynomial variables.
The verifier needs only the Python standard
library; its complete rational outputs accompany the article.

The separate numerical diagnostic evaluates the ordinary sums at
$90$ decimal digits, with $40$ or $64$ initial terms and fixed centered
asymptotic tail completions. For the four displayed cubic moments the
$40$-term residuals are between $3.4\times10^{-72}$ and
$1.6\times10^{-62}$; the $64$-term residuals are below $10^{-76}$.
The quartic second-moment relation has residual below $3.4\times10^{-72}$,
and the larger rational identity has residual below $1.3\times10^{-65}$.
These are numerical diagnostics, not interval certificates. The proofs
are the finite identities, the analytic continuation argument, and
exact rational elimination.

Several precise continuations remain:
\begin{enumerate}
\item Determine the smallest useful formal arithmetic space for a
      general multiset of even orders, and find a bound for the moment
      prefix needed to recover it. The examples above have dimensions
      six and eight, far below the unrestricted weight-based list.
\item Find further low-degree combinations of mixed moments whose
      largest-weight anchors cancel, as in
      \eqref{mtp:short-cubic-formula} and \eqref{mtp:short-quartic-formula}.
      A systematic classification could produce ordinary convergent
      identities with only single zeta values and their products.
\item Extend the boundary polynomial to common noninteger shifts.
      The shifted lattice no longer annihilates every nonnegative even
      power, so the finite endpoint correction must be retained.
\item Treat odd tail orders or spectral-order derivatives. Their
      centered parity and logarithmic asymptotic terms differ from the
      hypotheses proved here; simple polarization of the present
      theorem does not supply those cases.
\item Identify minimal functional multiple-polylogarithm depth under
      known functional relations. The bound $d-1$ is constructive,
      and makes no assertion that a smaller expression is impossible.
\end{enumerate}

% END INPUT


% BEGIN sections/06_independent_shifts.tex
% Standalone article section for the parent manuscript.
% Requires amsmath, amssymb, amsthm, hyperref and ordinary theorem/corollary/remark.
% New macros: \gsZ, \gsJ, \gsB, \gsV, \gsQ, \gsE, \gsT, \gsL.
% Bibliography keys: ShpotParis, Connon, DLMF, RepoTailKernels,
% RepoBilinear.  Proposed bibliography text is supplied separately.
\newcommand{\gsZ}{\mathcal Z}
\newcommand{\gsJ}{\mathcal J}
\newcommand{\gsB}{\mathcal B}
\newcommand{\gsV}{\mathcal V}
\newcommand{\gsQ}{\mathcal Q}
\newcommand{\gsE}{\mathcal E}
\newcommand{\gsT}{\mathcal T}
\newcommand{\gsL}{\mathcal L}

\section{Independent nonperiodic shifts: a completed product generator}
\label{gs:section}

The ordinary integral
\[
 \int_0^1\gamma_m^{(r)}(a+x)\gamma_n^{(t)}(b+x)\,dx,
 \qquad \Re a,\Re b>0,
\]
has two independent argument shifts.  It is a different object from a
periodically translated product.  Question R8 of \emph{Tail Kernels and
Boundary Products} asks for an order-generating formula, especially its
antisymmetric first-argument-derivative component
\cite{RepoTailKernels}.  This section supplies such a formula in an
explicit function class: finitely many Gauss hypergeometric terms,
holomorphic Gamma divided differences, and two geometrically convergent
single Hurwitz-zeta series.  It also gives a smaller function class for
all products of ordinary polygamma functions.

The mechanism extends the classical product-integral method of
Shpot and Paris \cite{ShpotParis}.  Their coincident shift is the
case $a=b=1$ before the spectral poles are removed.  The logarithmic
series occurring in digamma-square integrals is also classical; see
Connon \cite{Connon}.  The additions established here are the
independent-shift completion as a \emph{joint} holomorphic germ, its
explicit antisymmetric derivative, and the finite closure of all
ordinary argument derivatives in one logarithmic Hurwitz family.
These are claims relative to the inspected project sources; historical
priority for every specialization is not asserted.

\subsection{Unfolding the two-dimensional lattice}

For $\Re A,\Re B>0$ and $\Re(s+t)>1$, define the elementary
incomplete-beta kernel
\begin{align}
 \gsB_{s,t}(A,B)
 &=\int_0^\infty\frac{dy}{(y+A)^s(y+B)^t}\notag\\
 &=\frac{B^{1-s-t}}{s+t-1}\,
 {}_2F_1\left(s,s+t-1;s+t;1-\frac AB\right).
 \label{gs:beta-kernel}
\end{align}
The first expression fixes the branches of the second.  Thus the
formula is regular when $A=B$, with value
$A^{1-s-t}/(s+t-1)$.

\begin{theorem}[Exact diagonal unfolding]\label{gs:unfolding}
For $\Re s,\Re t>1$ and $\Re a,\Re b>0$,
\begin{equation}\label{gs:unfolded}
 \int_0^1\zeta(s,a+x)\zeta(t,b+x)\,dx
 =\sum_{k=0}^\infty\gsB_{s,t}(a+k,b)
  +\sum_{k=1}^\infty\gsB_{s,t}(a,b+k).
\end{equation}
Both sums converge absolutely and locally normally in this domain.
\end{theorem}

\begin{proof}
Expand both Hurwitz functions by their absolutely convergent Dirichlet
series.  On the lattice half-plane $m\ge n\ge0$, put $m=n+k$.
For each fixed $k\ge0$, the intervals $n+x$, $0\le x\le1$, tile the
positive real axis; consequently
\[
 \sum_{n=0}^\infty\int_0^1
 (x+n+k+a)^{-s}(x+n+b)^{-t}\,dx
 =\gsB_{s,t}(a+k,b).
\]
The complementary half-plane $n>m\ge0$ gives the second sum, starting
at $k=1$.  This starting index is essential: the lattice diagonal is
counted once.  Absolute convergence follows by replacing the complex
shifts with lower bounds for their real parts and applying the same
argument to positive integrands.  The bounds are uniform on compact
subsets, proving normal convergence.
\end{proof}

Choose integers $K\ge0$ and $L\ge1$, and write
\begin{equation}\label{gs:cd}
 c=K+a-b,\qquad d=L+b-a.
\end{equation}
They may always be chosen so large that
\begin{equation}\label{gs:cutoff-condition}
 \Re c>|b|,\qquad \Re d>|a|.
\end{equation}
The choices can be kept fixed on any sufficiently small neighborhood
of given shifts.  Put
\begin{equation}\label{gs:finite-prefix}
 \Phi_{s,t}(a,b)=
 \sum_{k=0}^{K-1}\gsB_{s,t}(a+k,b)
 +\sum_{k=1}^{L-1}\gsB_{s,t}(a,b+k),
\end{equation}
where an empty sum is zero.  A useful accelerated form of
\eqref{gs:unfolded} is
\begin{align}
 I(s,t;a,b)={}&\Phi_{s,t}(a,b)
 +B(1-t,s+t-1)\zeta(s+t-1,c)\notag\\
 &+B(1-s,s+t-1)\zeta(s+t-1,d)\notag\\
 &-\sum_{j=0}^\infty
 \frac{(-1)^j(s)_j b^{j+1-t}}{j!(j+1-t)}\zeta(s+j,c)\notag\\
 &-\sum_{j=0}^\infty
 \frac{(-1)^j(t)_j a^{j+1-s}}{j!(j+1-s)}\zeta(t+j,d),
 \label{gs:meromorphic-master}
\end{align}
where $I$ denotes the integral on the left of \eqref{gs:unfolded}.
Here $B$ without calligraphic type is Euler's beta function.
Identity \eqref{gs:meromorphic-master} is understood meromorphically;
its separate terms need not be regular at integral orders.

Indeed, for a fixed $q$ with $\Re q>|b|$, the change of variable
$y+b=z$ in \eqref{gs:beta-kernel}, followed by subtraction of the
interval $0<z<b$, gives
\begin{equation}\label{gs:strip-evaluation}
 \gsB_{s,t}(q+b,b)=q^{1-s-t}B(1-t,s+t-1)
 -\sum_{j=0}^\infty
 \frac{(-1)^j(s)_j b^{j+1-t}}{j!(j+1-t)}q^{-s-j}.
\end{equation}
This first holds where the beta integral and the subtracted interval
both converge, and then holds by meromorphic continuation.  Summing
over $q=c,c+1,\ldots$ proves \eqref{gs:meromorphic-master}; the other
tail follows by exchange of the two factors.  Alternatively one may
first subtract the $j=0$ term and continue the resulting normally
convergent remainder.  In either derivation, the condition
\eqref{gs:cutoff-condition} gives a geometric majorant for the
$j$-series on compact sets avoiding the displayed poles.

\subsection{Removal of both spectral poles without an iterated limit}

Use the standard convention
\begin{equation}\label{gs:stieltjes-convention}
 \gsZ(u,z)=\zeta(1+u,z)-\frac1u
 =\sum_{m=0}^\infty\frac{(-1)^m\gamma_m(z)}{m!}u^m.
\end{equation}
For $\Re z>0$ this is entire in $u$.  The required ordinary product
generator is
\begin{equation}\label{gs:target-master}
 \gsJ(u,v;a,b)=\int_0^1\gsZ(u,a+x)\gsZ(v,b+x)\,dx.
\end{equation}
We now evaluate it; its definition alone is not the claimed result.
The Hurwitz recurrence and argument-derivative formulas used below
have the conventions of DLMF, \S25.11 \cite{DLMF}.

Define the universal Gamma functions
\[
 \gsQ(u,v)=\frac{\Gamma(1+u+v)}{\Gamma(1+u)\Gamma(1+v)},
 \qquad \gsE(w)=\frac{\pi w}{\sin\pi w}.
\]
The following quotient has a joint removable singularity:
\begin{equation}\label{gs:universal-quotient}
 \gsT(u,v)=
 \frac{u+v-\gsQ(u,v)\{u\gsE(v)+v\gsE(u)\}}
 {uv(u+v)}.
\end{equation}
For an explicitly holomorphic interpretation, write
\begin{align}
 \gsT(u,v)&=-\frac{\gsQ(u,v)-1}{uv}
 -\gsQ(u,v)\sum_{h=1}^\infty e_{2h}
       \frac{u^{2h-1}+v^{2h-1}}{u+v},\label{gs:universal-holomorphic}\\
 e_{2h}&=2(1-2^{1-2h})\zeta(2h).
\end{align}
Each quotient in the series is a polynomial; also
$\gsQ-1$ is divisible by $uv$ in the ring of convergent power
series.  In particular $\gsT(0,0)=-2\zeta(2)$.

The only infinite series in the completed answer are
\begin{equation}\label{gs:V}
 \gsV(u,v;c,z)=
 -\sum_{j=1}^\infty
 \frac{(-1)^j(1+u)_j z^{j-v}}{j!(j-v)}
                    \zeta(1+u+j,c),\qquad \Re c>|z|.
\end{equation}
They are holomorphic near $(u,v)=(0,0)$, locally normally in all
variables.  No term at $j=0$ is included.

\begin{theorem}[Completed independent-shift generator]
\label{gs:completed}
Under \eqref{gs:cd}--\eqref{gs:cutoff-condition}, the ordinary integral
\eqref{gs:target-master} is the following joint holomorphic germ:
\begin{align}
 \boxed{\begin{aligned}
 \gsJ(u,v;a,b)={}&\Phi_{1+u,1+v}(a,b)+\gsT(u,v)\\
 &+\frac{b^{-v}\gsZ(u,c)
       -\gsQ(u,v)\gsE(v)\gsZ(u+v,c)}{v}\\
 &+\frac{a^{-u}\gsZ(v,d)
       -\gsQ(u,v)\gsE(u)\gsZ(u+v,d)}{u}\\
 &+\gsV(u,v;c,b)+\gsV(v,u;d,a).
 \end{aligned}}
 \label{gs:completed-formula}
\end{align}
Every displayed divided difference has its unique holomorphic value
on its zero-denominator face.  All mixed spectral and argument
derivatives may be taken term by term in the two $\gsV$ series.
The result is independent of the permissible integers $K,L$.
\end{theorem}

\begin{proof}
The unit-interval integral of one Hurwitz function is
\[
 \int_0^1\zeta(s,a+x)\,dx=\frac{a^{1-s}}{s-1}.
\]
Consequently, with $s=1+u$, $t=1+v$,
\begin{equation}\label{gs:pole-removal}
 \gsJ(u,v;a,b)=I(1+u,1+v;a,b)
             -\frac{a^{-u}+b^{-v}-1}{uv}.
\end{equation}
In \eqref{gs:meromorphic-master}, isolate $j=0$ in each series and
use
\[
 B(-v,1+u+v)=-\frac{\gsQ(u,v)\gsE(v)}v,
 \qquad
 B(-u,1+u+v)=-\frac{\gsQ(u,v)\gsE(u)}u.
\]
Replacing each Hurwitz function with its pole plus \eqref{gs:stieltjes-convention}
leaves precisely the two divided differences in
\eqref{gs:completed-formula}.  The remaining rational pole terms
combine to \eqref{gs:universal-quotient}.  Formula
\eqref{gs:universal-holomorphic} proves that this last apparent
singularity is jointly removable, including on $u+v=0$.
The numerator of the first divided difference vanishes identically
on $v=0$; that of the second vanishes on $u=0$.

For normal convergence, on a compact subset of
$\Re c>|z|$ choose $\rho<1$ with $|z|/\Re c\le\rho$.
For large $j$, the defining Hurwitz series gives a bound of the form
$C(\Re c)^{-j}$ times a fixed polynomial in $j$; the Pochhammer
ratio $(1+u)_j/j!$ also has polynomial growth uniformly for $u$ in
a compact neighborhood of zero.  Thus the summands in
\eqref{gs:V} are bounded by $C'j^M\rho^j$.  Cauchy's estimates
on a slightly larger compact set give the same conclusion after any
fixed mixed differentiation.  The finite prefix is holomorphic
because $\Re(2+u+v)>1$ near the origin.
The meromorphic identity therefore extends to a holomorphic equality.
Its left side is independent of $K,L$, proving the final assertion.
\end{proof}

\begin{corollary}[Every independent-shift Stieltjes product]
\label{gs:all-indices}
For all $m,n,r,t\ge0$ and $\Re a,\Re b>0$,
\begin{equation}\label{gs:all-indices-formula}
 \int_0^1\gamma_m^{(r)}(a+x)\gamma_n^{(t)}(b+x)\,dx
 =(-1)^{m+n}
 \left.\partial_a^r\partial_b^t\partial_u^m\partial_v^n
                  \gsJ(u,v;a,b)\right|_{u=v=0},
\end{equation}
where the right side is evaluated by the explicit convergent formula
\eqref{gs:completed-formula}.  All constants of integration and all
spectral regularizations are fixed by this formula.
\end{corollary}

This is a convergent reduction, not a claim of finite reduction to
ordinary zeta values at arbitrary shifts.  To obtain a prescribed
coefficient, expand the finite Gamma germ through the required degree,
use generalized Stieltjes constants at $c,d$, and extract that degree
in each geometrically convergent series.  The resulting coefficients
involve finite harmonic sums from $(1+u)_j$, powers of $\log a$ and
$\log b$, and Hurwitz order derivatives.  Each requested coefficient
therefore has an explicit, rapidly convergent single-sum expression.
More explicitly, the entire generalized-harmonic contribution is
encoded by the finite-product identity
\[
 \frac{(1+u)_j}{j!}
 =\exp\left\{\sum_{k=1}^\infty
        \frac{(-1)^{k+1}}{k}H_j^{(k)}u^k\right\},
 \qquad H_j^{(k)}=\sum_{h=1}^j h^{-k},\quad |u|<1.
\]

\subsection{The finite terms are polylogarithmic at integral orders}

There is no need to introduce parameter derivatives of a hypergeometric
function as additional undefined constants.  For positive $A>B$, put
$q=A-B$ and $z=q/A$.  Equation \eqref{gs:beta-kernel} becomes
\begin{equation}\label{gs:polylog-kernel}
 \gsB_{1+u,1+v}(A,B)
 =q^{-1-u-v}\int_0^z
       \frac{\exp\{(u+v)\log y-v\log(1-y)\}}{1-y}\,dy.
\end{equation}
Every fixed Taylor coefficient is a finite linear combination of
harmonic polylogarithms on the alphabet $\{0,1\}$ and logarithms;
this follows by the shuffle product for powers of $\log y$ and
$\log(1-y)$ followed by one integration with kernel $1/(1-y)$.
The coefficient of total degree $N$ has weight at most $N+1$ before
multiplication by powers of $\log q$.  Other positive $A,B$ follow
by exchange of the two factors, and the diagonal is elementary.
For example, writing $\ell=\log(A/B)$,
\begin{align}
 \gsB_{1,1}(A,B)&=\frac{\ell}{A-B},\notag\\
 \left.\partial_u\gsB_{1+u,1+v}(A,B)\right|_{0,0}
 &=-\frac{\ell\log A+\operatorname{Li}_2(1-B/A)}{A-B},\notag\\
 \left.\partial_v\gsB_{1+u,1+v}(A,B)\right|_{0,0}
 &=-\frac{\ell\log A+\operatorname{Li}_2(1-B/A)-\ell^2/2}{A-B}.
 \label{gs:dilog-prefix}
\end{align}
Thus the finite prefix contributes explicitly reducible polylogarithms,
whereas the independent-shift information in general remains in the
two single zeta series.

\subsection{An explicit first Stieltjes coefficient}

The first spectral coefficient already combines generalized Stieltjes
constants, harmonic numbers, dilogarithms, and Hurwitz order
derivatives.  Define
\[
 D_1(A,B)=-\frac{\log(A/B)\log A+
                   \operatorname{Li}_2(1-B/A)}{A-B},
 \qquad D_1(A,A)=-\frac{1+\log A}{A},
\]
and let $\Phi^{[1]}(a,b)$ be the finite prefix
\eqref{gs:finite-prefix} with every $\gsB$ replaced by $D_1$.
These formulas are real and continuous for all $A,B>0$.

\begin{corollary}[A harmonic--Stieltjes product identity]
\label{gs:gamma1-digamma}
For $a,b>0$, choose $K,L$ satisfying
\eqref{gs:cutoff-condition} and use \eqref{gs:cd}.  Then
\begin{align}
 \int_0^1\gamma_1(a+x)\psi(b+x)\,dx
 ={}&\Phi^{[1]}(a,b)+\zeta(3)-\gamma_2(c)-\frac12\gamma_2(d)
       \notag\\
 &+\log b\,\gamma_1(c)+\zeta(2)\psi(c)
   +\left\{\zeta(2)-\frac12\log^2a\right\}\psi(d)\notag\\
 &-\sum_{j=1}^\infty\frac{(-b)^j}{j}
       \{H_j\zeta(j+1,c)+\zeta'(j+1,c)\}\notag\\
 &-\sum_{j=1}^\infty\frac{(-a)^j}{j}
          \left(\frac1j-\log a\right)\zeta(j+1,d).
 \label{gs:gamma1-digamma-formula}
\end{align}
Here $H_j=\sum_{k=1}^j1/k$ and the prime on $\zeta$ denotes its
first spectral derivative.  Both infinite sums converge geometrically.
\end{corollary}

\begin{proof}
The left side is $\partial_u\gsJ(0,0;a,b)$ under the convention
\eqref{gs:stieltjes-convention}.  The universal term contributes
$\partial_u\gsT(0,0)=\zeta(3)$, and
$\partial_u(1+u)_j|_{u=0}=j!H_j$ accounts for the harmonic number.
On the face $v=0$, the first divided difference in
\eqref{gs:completed-formula} equals
\[
 -\{\log b+\psi(1+u)+\gamma\}\gsZ(u,c)
 -\partial_u\gsZ(u,c).
\]
Its derivative at zero is
$\log b\,\gamma_1(c)+\zeta(2)\psi(c)-\gamma_2(c)$.
The derivative of the second divided difference is
$\{\zeta(2)-\tfrac12\log^2a\}\psi(d)-\tfrac12\gamma_2(d)$.
The finite prefix follows from \eqref{gs:dilog-prefix}; differentiating
the two normally convergent $\gsV$ series gives the displayed sums.
\end{proof}

\subsection{A compact logarithmic Hurwitz family}

For an integer $p\ge1$, $c>0$, $z\ge0$, set
\begin{equation}\label{gs:L-definition}
 \gsL_p(c,z)=\sum_{n=0}^\infty
       \frac{\log(1+z/(n+c))}{(n+c)^p}.
\end{equation}
This is an ordinary absolutely convergent series; it extends
holomorphically to a neighborhood of the positive real parameter
domain.  When $|z|<\Re c$, it has the geometric zeta expansion
\begin{equation}\label{gs:L-zeta}
 \gsL_p(c,z)=-\sum_{j=1}^\infty
                    \frac{(-z)^j}{j}\zeta(p+j,c).
\end{equation}
It is convenient notation for an explicit convergent series, not a
claim that this logarithmic zeta family is historically new.

Put $R_p=\partial_z\gsL_p$.  Direct partial fractions give the
terminating rules
\begin{align}
 R_1(c,z)&=\frac{\psi(c+z)-\psi(c)}z,\label{gs:R1}\\
 R_p(c,z)&=\frac{\zeta(p,c)-R_{p-1}(c,z)}z\quad(p\ge2),\label{gs:Rp}\\
 \partial_c\gsL_p(c,z)
 &=-p\gsL_{p+1}(c,z)+R_p(c,z)-\zeta(p+1,c).
 \label{gs:L-derivative}
\end{align}
All apparent singularities at $z=0$ are removable.  In particular,
\begin{equation}\label{gs:Gamma-primitive}
 \gsL_1(c,z)=\int_0^z\frac{\psi(c+y)-\psi(c)}y\,dy,
\end{equation}
with its additive constant fixed by $\gsL_1(c,0)=0$.
This identifies the precise Gamma primitive needed by the independent
shifts.  Its derivative is the divided difference of the logarithmic
derivative of $\Gamma$.

For completeness, one two-parameter member generates the full family:
\begin{equation}\label{gs:L-generator}
 \sum_{p=1}^\infty w^{p-1}\gsL_p(c,z)
 =\gsL_1(c-w,z+w)-\gsL_1(c-w,w),\qquad |w|<c,
\end{equation}
locally where all logarithms have their anchored branches.  Summing
the geometric series $\sum_{p\ge1}w^{p-1}(n+c)^{-p}$ proves this
identity.  The finite shift rule
\begin{equation}\label{gs:L-shift}
 \gsL_p(c+1,z)=\gsL_p(c,z)-c^{-p}\log(1+z/c)
\end{equation}
makes evaluation at large $z/c$ straightforward: shift $c$ until
\eqref{gs:L-zeta} converges geometrically, retaining only finitely
many elementary logarithms.

\subsection{An explicit two-shift digamma product}

Define the continuous divided difference
\[
 D(A,B)=\frac{\log(A/B)}{A-B},\qquad D(A,A)=\frac1A.
\]
For the moment suppose $a,b>0$ and $|a-b|<1$, and abbreviate
\[
 c=1+a-b,\qquad d=1+b-a.
\]

\begin{theorem}[Ordinary independent-shift digamma identity]
\label{gs:digamma}
\begin{align}
 \boxed{\begin{aligned}
 \int_0^1\psi(a+x)\psi(b+x)\,dx
 ={}&D(a,b)+\gamma_1(c)+\gamma_1(d)-2\zeta(2)\\
 &+\log b\,\psi(c)+\log a\,\psi(d)\\
 &+\gsL_1(c,b)+\gsL_1(d,a).
 \end{aligned}}
 \label{gs:digamma-formula}
\end{align}
For arbitrary $a,b>0$, choose $K\ge0,L\ge1$ so that the $c,d$ in
\eqref{gs:cd} are positive, and replace $D(a,b)$ by
\[
 \sum_{k=0}^{K-1}D(a+k,b)+\sum_{k=1}^{L-1}D(a,b+k).
\]
All other terms retain the form in \eqref{gs:digamma-formula} with
these $c,d$.
\end{theorem}

\begin{proof}
Set $u=v=0$ in \eqref{gs:completed-formula}.
The first divided difference becomes
$\gamma_1(c)-\log b\,\gamma_0(c)$, because both first derivatives
of $\gsQ\gsE$ vanish at the origin.  The second is the exchanged
expression.  Also $\gsV(0,0;c,b)=\gsL_1(c,b)$,
$\gamma_0=-\psi$, and $\gsT(0,0)=-2\zeta(2)$.
This proves the assertion under the geometric conditions.
The ordinary logarithmic series extends it to all stated positive
parameters.  Independence of the cutoffs can alternatively be checked
without analytic continuation: increasing $K$ by one adds
$D(K+a,b)$ to the prefix, while the shift formulas for $\gamma_1$,
$\psi$, and \eqref{gs:L-shift} subtract exactly that same number.
The argument for $L$ is identical.
\end{proof}

The diagonal is particularly transparent:
\begin{equation}\label{gs:diagonal-digamma}
 \int_a^{a+1}\psi(x)^2\,dx
 =\frac1a-2\gamma\log a+2\gamma_1-2\zeta(2)
                         +2\gsL_1(1,a),\qquad a>0.
\end{equation}
This specialization belongs to the established digamma-square and
logarithmic-series literature \cite{Connon}; it is included as a
normalization check, not relabelled as a new result.

\subsection{The antisymmetric first-derivative component}

Let
\begin{equation}\label{gs:W0}
 W_0(A,B)=
 \frac{2\log(A/B)-(A-B)(A^{-1}+B^{-1})}{(A-B)^2},
 \qquad W_0(A,A)=0.
\end{equation}
For $a,b>0$, $|a-b|<1$, keep $c=1+a-b,d=1+b-a$.

\begin{theorem}[Explicit antisymmetric derivative]
\label{gs:Wronskian}
The ordinary integral
\[
 W(a,b)=\int_0^1
 \{\psi(a+x)\psi'(b+x)-\psi'(a+x)\psi(b+x)\}\,dx
\]
has the exact evaluation
\begin{align}
 \boxed{\begin{aligned}
 W(a,b)={}&W_0(a,b)\\
 &+2\bigl\{\zeta'(2,d)-\zeta'(2,c)
       +\log a\,\zeta(2,d)-\log b\,\zeta(2,c)\\
 &\hspace{19mm}+\gsL_2(c,b)-\gsL_2(d,a)\bigr\}\\
 &+\frac{2\psi(c)-\psi(a+1)}b
       +\frac{\psi(b+1)-2\psi(d)}a.
 \end{aligned}}
 \label{gs:Wronskian-formula}
\end{align}
For arbitrary positive shifts, use \eqref{gs:cd}, replace $W_0(a,b)$
by its two finite-prefix sums, and replace $\psi(a+1),\psi(b+1)$
by $\psi(K+a),\psi(L+b)$ respectively.
\end{theorem}

\begin{proof}
Apply $\partial_b-\partial_a$ to \eqref{gs:digamma-formula}.
The finite divided difference gives $W_0$.
For $c=K+a-b$ one has
$(\partial_b-\partial_a)c=-2$, and the corresponding derivative
of $d$ is $2$.  Formula \eqref{gs:L-derivative} gives
\[
 (\partial_b-\partial_a)\gsL_1(c,b)
 =2\gsL_2(c,b)-R_1(c,b)+2\zeta(2,c),
\]
and the exchanged term has the opposite form.  The standard identity
\[
 \gamma_1'(z)=\zeta(2,z)+\zeta'(2,z)
\]
follows by differentiating the Hurwitz shift derivative at its spectral
pole.  It cancels the explicit $\zeta(2,c),\zeta(2,d)$ terms just
produced.  Finally use \eqref{gs:R1}, $c+b=K+a$, and $d+a=L+b$.
Every differentiation is justified by normal convergence of the
ordinary logarithmic series on compact positive parameter sets.
\end{proof}

The symmetric companion is the elementary endpoint identity
\begin{equation}\label{gs:symmetric}
 \int_0^1
 \{\psi'(a+x)\psi(b+x)+\psi(a+x)\psi'(b+x)\}\,dx
 =\frac{\psi(a)}b+\frac{\psi(b)}a+\frac1{ab}.
\end{equation}
Combining \eqref{gs:Wronskian-formula} and \eqref{gs:symmetric}
therefore evaluates each ordered first-derivative product separately.
The antisymmetric formula contains information not supplied by
integration by parts alone.

\subsection{Closure under every ordinary argument derivative}

Write $q=r+t$.  In differentiating \eqref{gs:digamma-formula}, the
only non-elementary series that must remain after $q$ derivatives are
\begin{equation}\label{gs:highest-L}
 q!\{(-1)^r\gsL_{q+1}(c,b)+(-1)^t\gsL_{q+1}(d,a)\}.
\end{equation}
All other terms reduce finitely to polygamma values, rational
functions, logarithms, and Hurwitz-zeta order derivatives.  For
$q>0$, even the two generalized Stieltjes constants disappear:
\begin{equation}\label{gs:gamma1-derivatives}
 \gamma_1^{(q)}(z)=(-1)^{q+1}q!
        \{\zeta'(q+1,z)+H_q\zeta(q+1,z)\}.
\end{equation}
To prove \eqref{gs:highest-L}, note that the highest remaining
$\gsL$ term in a derivative with respect to $a$ of $\gsL_p(c,b)$
is $-p\gsL_{p+1}$, whereas a derivative with respect to $b$ has
highest term $p\gsL_{p+1}$.  The remaining $R_p$ terms terminate
by \eqref{gs:R1}--\eqref{gs:Rp}.  Induction gives the displayed
coefficient.  The other term is exchanged.  Equation
\eqref{gs:gamma1-derivatives} follows by taking the coefficient of
the linear spectral term in
$\partial_z^q\zeta(1+u,z)=(-1)^q(1+u)_q\zeta(q+1+u,z)$.
Thus the rules are a terminating reduction algorithm, rather than an
instruction to leave arbitrary derivatives unevaluated.

For example, a completely explicit trigamma-square identity is
\begin{align}
 \int_a^{a+1}\psi'(x)^2\,dx
 ={}&\frac1{3a^3}+4\zeta'(3)+\frac{4\zeta(2)}a
        +4\zeta(3)\log a\notag\\
 &-\frac{2\{\psi(1+a)+\gamma\}}{a^2}
        -4\gsL_3(1,a),\qquad a>0.
 \label{gs:trigamma-square}
\end{align}
It is obtained by differentiating the two independent shifts first
and only then setting $a=b$.  Merely differentiating
\eqref{gs:diagonal-digamma} twice would yield the sum of several
ordered derivative products, not this individual square.

\subsection{Diagonal limits and the inherited coincident finite parts}

There are two different diagonal operations.  Setting $a=b>0$
is an ordinary holomorphic specialization of
\eqref{gs:completed-formula}.  Sending the common shift to zero is
a boundary operation and requires removal of the local divergences.
For fixed indices $m,n,r,t$, write
\[
 P(x)=\gamma_m^{(r)}(x)\gamma_n^{(t)}(x).
\]
The exact recurrence
$\gamma_m(x)=x^{-1}\log^m x+\gamma_m(1+x)$ supplies its finite
singular expansion
\[
 P(x)=\sum_{\alpha\ge1,h\ge0}C_{\alpha,h}x^{-\alpha}\log^h x
           +O(1+|\log x|^M)
\]
near zero, with finitely many terms in the displayed singular sum.
The coefficients are explicit products of the Laurent singular
polynomials and Taylor derivatives at one.  For each monomial let
$A_{\alpha,h}(x)$ be an antiderivative with zero constant term in
its power-log expansion.  Then
\begin{equation}\label{gs:boundary-limit}
 \lim_{a\downarrow0}\left\{
 \int_0^1P(a+x)\,dx
       +\sum_{\alpha,h}C_{\alpha,h}A_{\alpha,h}(a)\right\}
 =\operatorname{FP}\int_0^1P(x)\,dx.
\end{equation}
Indeed, the shifted integral is $\int_a^{1+a}P$, and its difference
from $\int_a^1P$ tends to zero.  After integration, the singular
monomials contribute $-C_{\alpha,h}A_{\alpha,h}(a)$; the remaining
integral converges by the displayed local bound.  This proves the
boundary compatibility for every fixed order, in the original
coordinate $x$, without exchanging an unjustified spectral limit
with a finite-part operation.

For example, \eqref{gs:diagonal-digamma} and
\eqref{gs:trigamma-square} give the exact checks
\begin{align}
 \operatorname{FP}\int_0^1\psi(x)^2\,dx
     &=2\gamma_1-2\zeta(2),\notag\\
 \operatorname{FP}\int_0^1\psi(x)\psi'(x)\,dx
     &=\zeta(2),\notag\\
 \operatorname{FP}\int_0^1\psi'(x)^2\,dx
     &=2\zeta(3)+4\zeta'(3).
 \label{gs:boundary-checks}
\end{align}
These agree with the inherited coincident product generator.  They
are compatibility checks, not newly claimed coincident evaluations.
Unequal-rate collisions $a,b\to0$ are not covered by a claim that
one can omit the local subtractions; their path-dependent corrections
remain a separate question.

\subsection{Source reconciliation and further targets}

There is a concrete status correction to the incoming research list.
Question R7 of \emph{Tail Kernels and Boundary Products} asks for all
nonpositive-order continuations of the alternating square law.
The cutoff-jet section of the independently archived
\emph{Bilinear Tails and Determinants} already proves the requisite
normal subtractions, all cutoff contacts, their spectral jets, the
order-zero $1/4$ discrepancy, and the change under a translated cutoff
\cite{RepoBilinear}.  R7 should therefore be marked resolved by
that report, with its stated coordinate convention.  Its Gamma-square
specializations should not be counted again as new results here.

The completed generator resolves R8 in the explicit convergent function
class described above.  It does not claim that arbitrary independent
shifts reduce to finitely many ordinary zeta constants.  Three sharper
questions remain:
\begin{enumerate}
\item Determine special rational pairs $(a,b)$ for which the two
logarithmic Hurwitz series in \eqref{gs:Wronskian-formula} reduce
further to a finite collection of classical polylogarithms and Gamma
values.  The coincident and complementary-shift cases should be
separated from generic rational pairs.
\item Derive functional equations in the shift difference $a-b$
that relate finite descriptions of the prefix by polylogarithms
to the two geometric zeta series.  The exact $K,L$ invariance gives
necessary tests for any proposed relation.
\item Compute the local correction for $a=\lambda\epsilon$,
$b=\mu\epsilon$ with fixed positive $\lambda,\mu$ and compare its
constant term with the coincident finite part.  The generator makes
this a defined two-parameter collision problem, but the diagonal
compatibility theorem alone does not settle it.
\end{enumerate}

% END INPUT


% BEGIN sections/07_spectral_faces.tex
% Self-contained article sections.  Standard packages: amsmath, amssymb,
% amsthm, mathtools.  Standard theorem environments are used.
% Macros expected from parent: \Li, \dd, \ii, \Ree, \C, \N.
% No additional custom macros are required.  Labels use the prefix sfj:.

\section{Exact coefficients for every spectral face}
\label{sfj:sec:faces}

The newest incoming spectral-face report gives an all-order common-ray
formula and asks, at the end of its spectral section, for complete
collision coefficients and independently variable spectral orders
\cite{RepoBernoulli}.  We answer both questions.  A further
specialization evaluates every face with just one displaced mass by
an ordinary series of \emph{depth-two} polylogarithms.  The classical
Dirichlet integral, Gamma product, and polylogarithm differentiation
identities are ingredients, not novelty claims
\cite{Carlson,DLMF}.

\subsection{The inherited face normalization}

Let $k\ge2$, let $b_1,\ldots,b_k$ be real, and let
\[
 \Sigma_k=\{t_i\ge0:\ t_1+\cdots+t_k=1\}.
\]
The measure is $dt_1\cdots dt_{k-1}$, with the last coordinate
eliminated.  For a nonempty subset $T\subseteq\{1,\ldots,k\}$ put
\[
 t_T=\sum_{i\in T}t_i,\qquad
 b_T=\frac{\sum_{i\in T}t_i b_i}{t_T}.
\]
The perspective $t_TF(b_T)$ is filled by zero when $t_T=0$.  Define
\begin{align}
 \mathcal H_k(F;t)&=
 \sum_{\varnothing\ne T\subseteq\{1,\ldots,k\}}
       (-1)^{k-|T|}t_TF(b_T),\label{sfj:eq:face}\\
 I_k(F;b_1,\ldots,b_k)&=
 \int_{\Sigma_k}\mathcal H_k(F;t)\frac{dt}{t_1\cdots t_k}.
 \label{sfj:eq:I}
\end{align}
The notation $I_J$ denotes the same construction with the labels in a
subset $J$.  This is exactly the predecessor's face normalization; in
particular no Euclidean surface-area factor is inserted.

For $F\in C^1([\min b_i,\max b_i])$, the integral is ordinary and
absolutely convergent.  Indeed the perspective is Lipschitz on its
nonnegative coordinate cone, because
\[
 \partial_{t_i}\{t_TF(b_T)\}
      =F(b_T)+F'(b_T)(b_i-b_T)\quad(i\in T).
\]
Terms cancel in pairs when any coordinate is zero.  Consequently
$|\mathcal H_k(F;t)|\le C\min_i t_i$.  On a region where one
coordinate is maximal, it is at least $1/k$ and
$\min_i t_i\le(\prod_{i\ne i_{\max}}t_i)^{1/(k-1)}$.
This also proves that
\begin{equation}\label{sfj:eq:normal-face-domain}
 \int_{\Sigma_k}\mathcal H_k(F;t)\prod_i t_i^{\alpha_i-1}\,dt
 \quad\hbox{is holomorphic if}\quad
 \Ree\alpha_i>-\frac1{k-1}\ \hbox{for every }i,
\end{equation}
with locally normally convergent logarithmic derivatives.  These
elementary estimates justify every face integral used below.  Constants
and affine functions are annihilated by $\mathcal H_k$.

\subsection{A finite harmonic formula at every polynomial degree}

For $m\ge2$ and $\ell\ge0$ define
\begin{equation}\label{sfj:eq:complete-harmonic}
 \mathfrak h_\ell(m)
  =[v^\ell]\prod_{j=1}^{m-1}(1-v/j)^{-1}.
\end{equation}
Thus $\mathfrak h_\ell(m)$ is the complete homogeneous polynomial
of degree $\ell$ in $1,1/2,\ldots,1/(m-1)$.  In the usual generalized
harmonic notation, its first values are
\[
 1,\qquad H_{m-1},\qquad
 \frac{H_{m-1}^2+H_{m-1}^{(2)}}2,\qquad
 \frac{H_{m-1}^3+3H_{m-1}H_{m-1}^{(2)}+2H_{m-1}^{(3)}}6.
\]
It is also the weakly ordered harmonic sum
$H_{m-1}^{\star}(\{1\}^{\ell})$.

\begin{theorem}[Every monomial face integral]
\label{sfj:thm:monomial}
Let $k\ge2$ and $m\ge2$.  For a vector
$\boldsymbol h=(h_1,\ldots,h_k)\in\N_0^k$ let
$d(\boldsymbol h)=|\{i:h_i>0\}|$.  Then
\begin{equation}\label{sfj:eq:monomial}
 \boxed{\displaystyle
 I_k(B\mapsto B^m;\boldsymbol b)
  =m\!\sum_{h_1+\cdots+h_k=m}
   (-1)^{k-d}(k-d)!\,\mathfrak h_{k-d}(m)
                  \prod_{h_i>0}\frac{b_i^{h_i}}{h_i}.}
\end{equation}
The right side is a finite symmetric polynomial with rational
coefficients.  The formula extends to complex $b_i$ as a polynomial
identity.
\end{theorem}

\begin{proof}
The perspective is homogeneous of degree one.  Radial integration
therefore gives
\[
 I_k(F;\boldsymbol b)=
 \int_{(0,\infty)^k} e^{-\sum_i t_i}
            \mathcal H_k(F;t)\frac{dt_1\cdots dt_k}{\prod_i t_i};
\]
the radial factor is $\int_0^\infty e^{-\rho}\,d\rho=1$.
For a monomial of degree $m\ge2$, use
\[
 t_T b_T^m=\frac{(\sum_{i\in T}t_i b_i)^m}{t_T^{m-1}},\qquad
 t_T^{1-m}=\frac1{(m-2)!}\int_0^\infty
                      v^{m-2}e^{-vt_T}\,dv.
\]
Fix a monomial vector $\boldsymbol h$ with support $U$ and $d=|U|$.
The alternating sum over $T\supseteq U$ is
\[
 e^{-v\sum_{i\in U}t_i}
             \prod_{j\notin U}(e^{-vt_j}-1).
\]
Integrating the active and inactive coordinates respectively gives
\[
 \frac{\Gamma(h_i)}{(1+v)^{h_i}},\qquad
 \int_0^\infty e^{-t}\frac{e^{-vt}-1}{t}\,dt=-\log(1+v).
\]
The last identity follows by differentiating in $v$ and fixing the
value at zero.  Every resulting term is absolutely integrable: its
remaining majorant is a constant times
$v^{m-2}(1+v)^{-m}\log^{k-d}(1+v)$.  Thus the use of Fubini is
legitimate after the finite inclusion--exclusion has been combined.
We obtain
\begin{align*}
 I_k(B^m;\boldsymbol b)
  =m(m-1)\sum_{|\boldsymbol h|=m}
     (-1)^{k-d}\prod_{h_i>0}\frac{b_i^{h_i}}{h_i}
       \int_0^\infty
        \frac{v^{m-2}\log^{k-d}(1+v)}{(1+v)^m}\,dv.
\end{align*}
On setting $x=(1+v)^{-1}$, the integral becomes
\[
 \int_0^1(1-x)^{m-2}(-\log x)^{\ell}\,dx
       =\frac{\ell!}{m-1}\mathfrak h_\ell(m),\qquad \ell=k-d.
\]
To verify the last equality, expand the elementary beta integral
\[
 \int_0^1x^u(1-x)^{m-2}\,dx
 =\frac1{m-1}\prod_{j=1}^{m-1}(1+u/j)^{-1}
\]
at $u=0$.  Substitution proves the stated formula.
\end{proof}

At $m=2$, Theorem~\ref{sfj:thm:monomial} reduces to the predecessor's
factorial variance identity
\begin{equation}\label{sfj:eq:quadratic}
 I_k(B^2;\boldsymbol b)
   =(-1)^{k-1}(k-2)!\sum_{i<j}(b_i-b_j)^2.
\end{equation}
The new formula supplies every degree beyond that quadratic term.  For
example,
\begin{equation}\label{sfj:eq:three-polynomials}
 I_3(B^2;0,1,2)=6,\qquad I_3(B^3;0,1,2)=18,\qquad
 I_3(B^4;0,1,2)=\frac{97}{2},\qquad
 I_3(B^5;0,1,2)=\frac{245}{2}.
\end{equation}

\subsection{Complete convergent collision expansions}

\begin{corollary}[All collision coefficients]
\label{sfj:cor:collision}
Suppose $F$ is holomorphic on $|z-B|<R$ and let
$M=\max_i|\beta_i|$.  For $|\epsilon|M<R$,
\begin{equation}\label{sfj:eq:collision}
 I_k(F;B+\epsilon\beta_1,\ldots,B+\epsilon\beta_k)
  =\sum_{m=2}^\infty
       \frac{F^{(m)}(B)}{m!}\epsilon^m
                     I_k(B^m;\boldsymbol\beta),
\end{equation}
normally on compact subsets of this disk.  Every coefficient is given
by the finite harmonic formula \eqref{sfj:eq:monomial}.
\end{corollary}
\begin{proof}
The constant and affine Taylor terms vanish.  On each smaller closed
disk, the Taylor series and its first derivative converge uniformly.
The boundary estimate preceding \eqref{sfj:eq:normal-face-domain}
is a continuous $C^1$ estimate and permits termwise integration.
Cauchy estimates on an intermediate disk give a geometric majorant,
including the harmless linear factor in $m$ from the derivative.
\end{proof}

The Gamma example is especially explicit.  For $a>0$ and $B\ge0$ put
\begin{align}
 G_a(B)&=\log\Gamma(a+\ii\sqrt B)
          +\log\Gamma(a-\ii\sqrt B)-\log(2\pi),
                    \label{sfj:eq:gammaG}\\
 R_m(a;B)&=\sum_{n=0}^\infty\bigl((n+a)^2+B\bigr)^{-m}.
                    \label{sfj:eq:resolvent}
\end{align}
The Gamma logarithms are analytic logarithms continued from positive
real arguments.  The paired function is analytic as a function of
$B$ on the disk of radius $a^2+B_0$ about any $B_0\ge0$.  The
Gamma product, or differentiation of its normally convergent logarithmic
difference, gives
\begin{equation}\label{sfj:eq:gamma-derivatives}
 G_a^{(m)}(B)=(-1)^m(m-1)!R_m(a;B),\qquad m\ge1.
\end{equation}
Consequently
\begin{equation}\label{sfj:eq:gamma-collision}
 \boxed{\displaystyle
 I_k(G_a;B+\epsilon\boldsymbol\beta)
   =\sum_{m=2}^\infty\frac{(-1)^m\epsilon^m}{m}
                  R_m(a;B)I_k(B^m;\boldsymbol\beta),
 \quad |\epsilon|M<a^2+B.}
\end{equation}
At $B=0$ the resolvent is simply $\zeta(2m,a)$.  For instance,
\begin{equation}\label{sfj:eq:gamma-three-example}
 I_3(G_a;0,\epsilon,2\epsilon)
   =3\zeta(4,a)\epsilon^2-6\zeta(6,a)\epsilon^3
       +\frac{97}{8}\zeta(8,a)\epsilon^4
       -\frac{49}{2}\zeta(10,a)\epsilon^5+\cdots,
 \qquad 2|\epsilon|<a^2.
\end{equation}
The positive integer resolvents themselves have a finite polygamma
description.  For $B>0$,
\begin{equation}\label{sfj:eq:resolvent-digamma}
 R_1(a;B)=
 \frac{\psi(a+\ii\sqrt B)-\psi(a-\ii\sqrt B)}{2\ii\sqrt B},
 \qquad
 R_m(a;B)=\frac{(-1)^{m-1}}{(m-1)!}
                 \partial_B^{m-1}R_1(a;B).
\end{equation}
Their apparent singularities at $B=0$ are removable.

\subsection{All logarithmic face moments have a rational compiler}

For independent parameters $\alpha_1,\ldots,\alpha_k$ define
\[
 A=\sum_i\alpha_i,\qquad
 Q_k(\boldsymbol\alpha)=
       \frac{\Gamma(1+A)}{\prod_i\Gamma(1+\alpha_i)},\qquad
 K_{k,m}(\boldsymbol\alpha;\boldsymbol b)
   =Q_k(\boldsymbol\alpha)
       \int_{\Sigma_k}\mathcal H_k(B^m;t)
                                      \prod_i t_i^{\alpha_i-1}\,dt.
\]
For an increment $v$, let
$\Delta_v f(y)=(f(y+v)-f(y))/v$, with $\Delta_0f=f'$.

\begin{theorem}[Rational normalized face moments]
\label{sfj:thm:rational}
Let $m\ge2$ and $f_m(y)=(1+y)_{m-1}^{-1}$.  For
$U=\{i:h_i>0\}$, write $A_U=\sum_{i\in U}\alpha_i$.  Then
\begin{equation}\label{sfj:eq:rational}
 \boxed{\displaystyle
 K_{k,m}=m!\sum_{|\boldsymbol h|=m}
  \left(\prod_{i=1}^k\frac{b_i^{h_i}}{h_i!}\right)
  \left(\prod_{i\in U}(1+\alpha_i)_{h_i-1}\right)
  \left(\prod_{j\notin U}\Delta_{\alpha_j}\right)f_m(A_U).}
\end{equation}
In particular this Gamma-normalized integral is a rational function of
all the spectral parameters, holomorphic in a neighborhood of zero.
Its Taylor coefficients give all its logarithmically weighted face
moments by finite rational algebra.  Removing $Q_k$ introduces only
the finite zeta coefficients in the Taylor series of $\log\Gamma$.
\end{theorem}
\begin{proof}
Repeat the orthant proof of Theorem~\ref{sfj:thm:monomial} with powers
$t_i^{\alpha_i-1}$.  The radial factor is $\Gamma(1+A)$, and the
normalization leaves division by $\prod_i\Gamma(1+\alpha_i)$.
For $i\in U$ the coordinate integral contributes
$(1+\alpha_i)_{h_i-1}(1+v)^{-h_i-\alpha_i}$.  For $j\notin U$
it contributes
$((1+v)^{-\alpha_j}-1)/\alpha_j$.
After $x=(1+v)^{-1}$, the remaining integral is
\[
 \int_0^1(1-x)^{m-2}x^{A_U}
            \prod_{j\notin U}\frac{x^{\alpha_j}-1}{\alpha_j}\,dx.
\]
It equals
$(m-2)!\,(\prod_{j\notin U}\Delta_{\alpha_j})f_m(A_U)$ by the
ordinary beta integral.  This proves the formula first where the
coordinate integrations are absolute.  The face estimate and the
removable finite differences extend both sides holomorphically to
zero.  At zero, $f_m^{(\ell)}(0)=
(-1)^\ell\ell!\mathfrak h_\ell(m)/(m-1)!$, which recovers
\eqref{sfj:eq:monomial} independently.
\end{proof}

\section{A depth-two polylogarithm evaluation of every repeated-mass face}
\label{sfj:sec:polylog}

The general collision expansion has a stronger closed consequence when
all but one mass agree.  It holds at arbitrary nonnegative displacement,
beyond the disk used by the collision series.

\begin{lemma}[One-dimensional reduction of a repeated-mass face]
\label{sfj:lem:repeated}
Let $F\in C^2([B,B+\delta])$, where $\delta\ge0$.  For $k\ge2$,
\begin{equation}\label{sfj:eq:one-dimensional}
 \boxed{\displaystyle
 I_k(F;\underbrace{B,\ldots,B}_{k-1},B+\delta)
  =(-1)^{k-1}\int_0^1
   \frac{d}{dt}\left\{\frac{F(B+\delta t)-F(B)}t\right\}
                            [-\log(1-t)]^{k-1}\,dt.}
\end{equation}
The quotient is filled by its limiting value at zero; the integral is
ordinary and absolutely convergent.
\end{lemma}
\begin{proof}
After translating $B$ and scaling $\delta$, it suffices to check a
polynomial $F$.  Constants and affine terms vanish on both sides.
For $F(t)=t^m$, $m\ge2$, only the composition supported on the single
nonzero mass survives in \eqref{sfj:eq:monomial}; its value is
$(-1)^{k-1}(k-1)!\mathfrak h_{k-1}(m)$.  The right side gives
\[
 (-1)^{k-1}(m-1)\int_0^1
           t^{m-2}[-\log(1-t)]^{k-1}\,dt,
\]
which has exactly the same value by the beta coefficient calculation.
For general $C^2$ functions, approximate $F''$ uniformly by
polynomials and integrate twice, preserving $F(B)$ and $F'(B)$.
The face estimate gives continuity of the left side in the $C^1$
norm.  On the right,
\[
 \frac{d}{dt}\frac{F(B+\delta t)-F(B)}t
       =\delta^2\int_0^1 uF''(B+\delta tu)\,du,
\]
so uniform $C^2$ convergence and integrability of every fixed power of
$-\log(1-t)$ give continuity there as well.  This proves the lemma.
\end{proof}

We use the strict-index double polylogarithm
\begin{equation}\label{sfj:eq:double-polylog}
 \Li_{1,p}(z)=\sum_{m>n\ge1}\frac{z^m}{m n^p}
    =\sum_{m=2}^\infty\frac{H_{m-1}^{(p)}}m z^m,
     \qquad |z|<1.
\end{equation}
In particular
$\frac{d}{dz}\Li_{1,p}(z)=\Li_p(z)/(1-z)$.

\begin{theorem}[Every repeated-mass Gamma face has depth two]
\label{sfj:thm:gamma-polylog}
Let $a>0$, $B\ge0$, $\delta\ge0$, and $k\ge2$.  With $G_a$ as in
\eqref{sfj:eq:gammaG},
\begin{equation}\label{sfj:eq:gamma-polylog}
 \boxed{\displaystyle
 I_k(G_a;\underbrace{B,\ldots,B}_{k-1},B+\delta)
  =(-1)^{k-1}(k-1)!
     \sum_{n=0}^\infty
       \Li_{1,k-1}\!\left(
             \frac{\delta}{(n+a)^2+B+\delta}\right).}
\end{equation}
The sum is ordinary and absolutely convergent.  It converges normally
when $(a,B,\delta)$ ranges over a compact subset of
$\{a>0,B\ge0,\delta\ge0\}$.  Thus the dimension of the face does
not force polylogarithm depth greater than two in this parameter sector.
\end{theorem}

\begin{proof}
The Gamma product gives the normally convergent logarithmic difference
\begin{equation}\label{sfj:eq:gamma-difference}
 G_a(B+\delta t)-G_a(B)
       =-\sum_{n=0}^\infty
          \log\left(1+\frac{\delta t}{(n+a)^2+B}\right).
\end{equation}
For $u\ge0$, define
\[
 J_k(u)=\int_0^1
  \frac{\log(1+ut)-ut/(1+ut)}{t^2}
                            [-\log(1-t)]^{k-1}\,dt.
\]
Its integrand is nonnegative; the apparent singularity at $t=0$ is
removable.  Differentiation under the integral is justified on compact
$u$-intervals and gives
\[
 J_k'(u)=u\int_0^1\frac{[-\log(1-t)]^{k-1}}{(1+ut)^2}\,dt.
\]
Set $q=u/(1+u)$ and $y=1-t$.  Expanding $(1-qy)^{-2}$, with
nonnegative terms, yields
\[
 J_k'(u)
  =\frac{(k-1)!}{1+u}\Li_{k-1}\!\left(\frac{u}{1+u}\right).
\]
Since
\[
 \frac{d}{du}\Li_{1,k-1}\!\left(\frac{u}{1+u}\right)
      =\frac1{1+u}\Li_{k-1}\!\left(\frac{u}{1+u}\right)
\]
and both functions vanish at zero, we have
$J_k(u)=(k-1)!\Li_{1,k-1}(u/(1+u))$.
Insert \eqref{sfj:eq:gamma-difference} into
Lemma~\ref{sfj:lem:repeated}; positivity permits Tonelli's theorem,
and this scalar identity gives \eqref{sfj:eq:gamma-polylog}.

For the convergence assertion, all arguments lie in $[0,q_0]$, where
$q_0=\delta/(a^2+B+\delta)<1$.  From
$H_{m-1}^{(p)}/m\le1$ for $p\ge1$,
\[
 0\le\Li_{1,p}(q)\le\frac{q^2}{1-q}.
\]
The summands are therefore $O((n+a)^{-4})$, uniformly on the stated
compact parameter sets.  This proves finiteness as well as normal
convergence.  The case $\delta=0$ is immediate.
\end{proof}

For $k=2$ this is the previously known log-square covariance,
since $\Li_{1,1}(q)=\frac12\log^2(1-q)$.  For $k=3$, the double
polylogarithm reduces entirely to classical dilogarithms and
trilogarithms:
\begin{equation}\label{sfj:eq:Li12-classical}
 \Li_{1,2}(q)=2\Li_3(1-q)-2\zeta(3)
            -\log(1-q)\{\Li_2(1-q)+\zeta(2)\},\qquad 0<q<1.
\end{equation}
The formula follows by differentiating both sides and taking the
limit at $q=0$.  It provides an exact classical-polylogarithm lattice
evaluation of the first genuinely higher face.

For example, at $a=1$, $B=0$, and $\delta=3/10$,
\begin{equation}\label{sfj:eq:numeric-example}
 I_3(G_1;0,0,3/10)
   =2\sum_{n=1}^\infty\Li_{1,2}\!\left(\frac{3}{10n^2+3}\right)
   =0.0727899864111738847269555921654165064396\ldots.
\end{equation}
The decimal is a diagnostic value, not an interval enclosure.

\begin{remark}[A weighted extension]
The same proof applies when $G_a$ is replaced by the first spectral
jet $F_1(B)$ of $\sum_{n\ge0}W(n+a)((n+a)^2+B)^{-s}$ and
$\deg W\le2$.  Its second mass derivative is the ordinary series
$\sum W(n+a)((n+a)^2+B)^{-2}$.  More explicitly, with $x_n=n+a$,
\[
 F_1(B+\delta t)-F_1(B)-\delta tF_1'(B)
 =-\sum_{n=0}^{\infty}W(x_n)
   \left\{\log\left(1+\frac{\delta t}{x_n^2+B}\right)
                       -\frac{\delta t}{x_n^2+B}\right\}.
\]
This twice-anchored sum and its first two $t$ derivatives converge
normally on compact displacement intervals, since their summands are
$O(x_n^{\deg W-4})$.  Equality follows from the common second
derivative and the two zero initial values at $t=0$.  Every face
annihilates the two displayed affine terms.  The right side of
\eqref{sfj:eq:gamma-polylog}
then acquires the factor $W(n+a)$.  Its convergence follows from
$W(n+a)=O(n^2)$.  At higher polynomial degree, additional subtractions
are necessary and are not silently included in this statement.
\end{remark}

\section{Independent orders: the exact polar obstruction and mixed jets}
\label{sfj:sec:independent}

The common-ray family can be holomorphic at zero for every fixed ray
even when the joint family is not a holomorphic germ.  This is a
material distinction in extending the predecessor's theorem.

\subsection{A complete holomorphy criterion}

Fix $a>0$, $r\ge2$, masses $B_i\ge0$, and a polynomial
$W(x)=\sum_{d=0}^Dw_dx^d$.  Put $x_n=n+a$, $A=\sum_i\alpha_i$,
and initially, for $\Ree A>(D+1)/2$, set
\begin{equation}\label{sfj:eq:joint}
 \mathscr Z_W(\boldsymbol\alpha)
    =\sum_{n=0}^\infty W(x_n)
                  \prod_{i=1}^r(x_n^2+B_i)^{-\alpha_i}.
\end{equation}
All quadratic factors are positive, and their powers use real
logarithms.  Define the finite polynomials
\begin{equation}\label{sfj:eq:polar-poly}
 C_h(\boldsymbol\alpha)
  =[v^h]\prod_{i=1}^r(1+B_iv)^{-\alpha_i},\qquad
 P_W(\boldsymbol\alpha)=\sum_{h\ge1}w_{2h-1}C_h(\boldsymbol\alpha).
\end{equation}
Only $h\le(D+1)/2$ occurs in $P_W$.

\begin{theorem}[The complete joint polar part]
\label{sfj:thm:pole}
Near $\boldsymbol\alpha=0$ there is a holomorphic function
$\mathscr H_W$ such that
\begin{equation}\label{sfj:eq:pole}
 \boxed{\displaystyle
 \mathscr Z_W(\boldsymbol\alpha)
     =\frac{P_W(\boldsymbol\alpha)}{2A}
                         +\mathscr H_W(\boldsymbol\alpha),\qquad
 \mathscr H_W(0)=\sum_{d=0}^D w_d\zeta(-d,a).}
\end{equation}
The quotient is a prescribed polar subtraction; its part divisible
by $A$ is of course holomorphic.  The joint spectrum is holomorphic
at zero if and only if either $W$ is even or all the masses $B_i$
are equal.
\end{theorem}

\begin{proof}
Choose a finite head so that $B_i/x_n^2<\rho<1$ on the tail.  A
finite binomial expansion with an analytic remainder gives
\[
 \prod_i(x_n^2+B_i)^{-\alpha_i}
     =x_n^{-2A}\sum_{h=0}^{K}C_h(\boldsymbol\alpha)x_n^{-2h}
                      +\text{remainder}.
\]
Take $K$ sufficiently large that the weighted remainder and all
required parameter derivatives converge normally near zero; any
$K$ with $2K+2>D+1$ leaves a strict convergence margin after a
small neighborhood is chosen.  The removed terms are finite sums of
$w_dC_h(\boldsymbol\alpha)\zeta(2A+2h-d,a+N)$, plus the retained
head.  Only $2h-d=1$ gives a pole through the origin, and its
residue is exactly $w_dC_h/(2A)$.  Subtracting these terms proves
\eqref{sfj:eq:pole}.  At zero, $C_0=1$ and $C_h(0)=0$ for $h\ge1$;
the finite head and tail then give the displayed value of
$\mathscr H_W(0)$.

If $W$ is even, the polar polynomial vanishes.  If every mass is
$B$, then $C_h=(-1)^h(A)_hB^h/h!$, which is divisible by $A$ for
$h\ge1$.  These prove sufficiency.  Conversely, suppose $w_{2H-1}$
is the highest nonzero odd coefficient and $B_i\ne B_j$.  Restrict
to $\alpha_i=t$, $\alpha_j=-t$, and all other parameters zero.
Then $A=0$ and
\[
 \prod_\ell(1+B_\ell v)^{-\alpha_\ell}
  =\exp\{t[\log(1+B_jv)-\log(1+B_iv)]\}.
\]
The coefficient of $t^H$ in $C_H$ is
$(B_j-B_i)^H/H!$.  Lower $C_h$ have degree at most $h<H$ and
cannot cancel it.  Thus $P_W$ is not divisible by $A$, proving
necessity.
\end{proof}

For a fixed vector $\boldsymbol\lambda$ with
$L=\sum_i\lambda_i\ne0$, every straight ray has the finite limit
\begin{equation}\label{sfj:eq:direction-limit}
 \lim_{t\to0}\mathscr Z_W(t\boldsymbol\lambda)
 =\sum_dw_d\zeta(-d,a)+\frac1L\sum_i\lambda_i p_W(B_i),\qquad
 p_W(B)=\sum_{h\ge1}\frac{(-1)^hw_{2h-1}}{2h}B^h.
\end{equation}
This follows from the linear coefficient
$C_h(t\boldsymbol\lambda)=
t(-1)^h\sum_i\lambda_iB_i^h/h+O(t^2)$.
For example, $a=1$, $W=x$, and $(B_1,B_2)=(0,1)$ give
\begin{equation}\label{sfj:eq:direction-counterexample}
 \lim_{t\to0}\mathscr Z_x(t,t)=-\frac13,
 \qquad
 \lim_{t\to0}\mathscr Z_x(t,2t)=-\frac5{12}.
\end{equation}

Even agreement of all straight-ray limits does not suffice.  For
$W=x^3+x/2$ and $(B_1,B_2)=(0,1)$, the polar term is
$\alpha_2^2/(4(\alpha_1+\alpha_2))$.  Every non-tangential straight
ray has the same zero-order limit, but along
$\alpha_2=t$, $\alpha_1=-t+t^3$ this term is $1/(4t)$.
The joint germ is therefore singular.  These are obstructions to an
unqualified extension, not errors in the correctly scoped common-ray
theorem of the predecessor.

\begin{corollary}[Periodic mean cancellation]
If a bounded periodic weight $c_n$ is inserted in every sum, with
mean $\bar c$, the polar term in \eqref{sfj:eq:pole} is multiplied
by $\bar c$.  The joint germ is holomorphic exactly when
$\bar c=0$, or $W$ is even, or all masses agree.
\end{corollary}
\begin{proof}
For a period $q$, the auxiliary Dirichlet series is
$q^{-z}\sum_{v=0}^{q-1}c_v\zeta(z,(a+v)/q)$ and its residue at
$z=1$ is $\bar c$.  All other parts of the proof are unchanged.
\end{proof}

\subsection{The holomorphic mixed defect after an explicit subtraction}

Write
\[
 E_W(s;B)=\sum_{n\ge0}W(x_n)(x_n^2+B)^{-s},\qquad
 F_j(B)=\left.\partial_s^jE_W(s;B)\right|_{s=0}.
\]
The one-mass functions are holomorphic at zero and
$F_0(B)=\sum_dw_d\zeta(-d,a)+p_W(B)$.  Define
\begin{equation}\label{sfj:eq:joint-defect}
 \mathscr D_W(\boldsymbol\alpha)
   =\frac1A\sum_i\alpha_iE_W(A;B_i)
                           -\mathscr Z_W(\boldsymbol\alpha).
\end{equation}
Let $p_h=(-1)^hw_{2h-1}/(2h)$, and set
\begin{equation}\label{sfj:eq:joint-local}
 \mathscr L_W(\boldsymbol\alpha)
  =\sum_{h\ge1}p_h\left\{
    \frac{\sum_i\alpha_iB_i^h}{A}
    -\frac{h!}{(A)_h}
       \sum_{h_1+\cdots+h_r=h}
          \prod_i\frac{(\alpha_i)_{h_i}B_i^{h_i}}{h_i!}
                         \right\}.
\end{equation}
This is a finite rational function.  Its $h=1$ term is identically
zero.  In several variables its apparent origin singularities should
\emph{not} in general be filled by a common-ray value.  For instance,
\begin{equation}\label{sfj:eq:joint-local-cubic}
 \mathscr L_{x^3}(\boldsymbol\alpha)
   =\frac{\sum_{i<j}\alpha_i\alpha_j(B_i-B_j)^2}{4A(1+A)}.
\end{equation}

\begin{theorem}[Independent-order face decomposition]
\label{sfj:thm:joint-faces}
Put
\[
 \mathcal G_W(u;B)=\frac{E_W(u;B)-F_0(B)}u
       =\sum_{j\ge1}\frac{F_j(B)}{j!}u^{j-1},
 \qquad
 Q_J(\boldsymbol\alpha)=
 \frac{\Gamma(1+\sum_{i\in J}\alpha_i)}
                      {\prod_{i\in J}\Gamma(1+\alpha_i)}.
\]
The difference $\widetilde{\mathscr D}_W=
\mathscr D_W-\mathscr L_W$ is holomorphic near zero and equals
\begin{equation}\label{sfj:eq:joint-faces}
 \boxed{\displaystyle
 \widetilde{\mathscr D}_W(\boldsymbol\alpha)
   =-\sum_{\substack{J\subseteq\{1,\ldots,r\}\\|J|\ge2}}
       \left(\prod_{i\in J}\alpha_i\right)Q_J(\boldsymbol\alpha)
       \int_{\Sigma_J}\mathcal H_J(\mathcal G_W(A;\cdot);t)
                              \prod_{i\in J}t_i^{\alpha_i-1}\,dt.}
\end{equation}
A sufficient neighborhood for all displayed integrals and zero-point
germs is $\max_i|\alpha_i|<1/(4r)$.  The formula supplies ordinary,
absolutely convergent integrals and locally normal derivatives there.
\end{theorem}

\begin{proof}
For initially positive real parts, the usual Dirichlet parameterization
of the product in \eqref{sfj:eq:joint} has total order $A$.
Finite subset inversion decomposes the perspective into its faces.
Under the Dirichlet measure with parameters $\alpha_i$, the mass
$t_J$ has mean $A_J/A$, where $A_J=\sum_{i\in J}\alpha_i$, and
is independent of the normalized coordinates within $J$.  Thus the
contribution of face $J$ is
\[
 \frac{A_J}{A}\frac{\Gamma(A_J)}{\prod_{i\in J}\Gamma(\alpha_i)}
 \int_{\Sigma_J}\mathcal H_J(E_W(A;\cdot);t)
                                         \prod_{i\in J}t_i^{\alpha_i-1}\,dt
 =\frac{\prod_{i\in J}\alpha_i}{A}Q_J
                    \int_{\Sigma_J}\mathcal H_J(E_W(A;\cdot);t)
                                         \prod_{i\in J}t_i^{\alpha_i-1}\,dt.
\]
The singleton terms give the weighted endpoint average in
\eqref{sfj:eq:joint-defect}.  Splitting $E_W(A;B)=F_0(B)+
A\mathcal G_W(A;B)$ yields \eqref{sfj:eq:joint-faces}; ordinary
Dirichlet polynomial moments give exactly \eqref{sfj:eq:joint-local}
for the $F_0$ part.  The identity continues meromorphically from its
common absolute-convergence domain.  In the stated polydisk,
$|A|<1/4$, each one-mass zero germ is holomorphic, and
\eqref{sfj:eq:normal-face-domain} permits every face integral and
its derivatives.  This proves the claimed joint holomorphy.
\end{proof}

For even $W$, the local term is zero and both the joint spectrum and
the joint defect are holomorphic.  More generally, the \emph{raw
defect} $\mathscr D_W$ is holomorphic at zero precisely when the odd
part of $W$ has degree at most one, or all masses agree.  With periodic
weights, mean zero is an additional sufficient and necessary alternative.
Indeed its polar numerator is the negative of the degree-at-least-two
part of $P_W$; the highest-degree restriction used in
Theorem~\ref{sfj:thm:pole} proves necessity.  This criterion explains
why the spectrum and its defect have different analytic boundaries.

\subsection{A particularly simple formula for distinct mixed derivatives}

For a set $J$ put
\[
 \Pi_J(\boldsymbol\alpha;t)=
 \exp\left\{\sum_{i\in J}\alpha_i\log t_i+
  \sum_{n\ge2}\frac{(-1)^n\zeta(n)}n
    \left[\left(\sum_{i\in J}\alpha_i\right)^n-
                             \sum_{i\in J}\alpha_i^n\right]\right\}.
\]
This is $Q_J\prod_{i\in J}t_i^{\alpha_i}$.  If $\boldsymbol\nu$
is a nonnegative multiindex, extraction from
\eqref{sfj:eq:joint-faces} gives the finite identity
\begin{align}\label{sfj:eq:all-mixed}
 \partial^{\boldsymbol\nu}\widetilde{\mathscr D}_W(0)
 ={}&-\boldsymbol\nu!\!
 \sum_{\substack{J\subseteq\operatorname{supp}\boldsymbol\nu\\|J|\ge2}}
 \sum_{j=1}^{|\boldsymbol\nu|-|J|+1}\frac1{j!}
 \int_{\Sigma_J}\mathcal H_J(F_j;t)
       [\boldsymbol\alpha^{\boldsymbol\nu-\boldsymbol1_J}]
        \{A^{j-1}\Pi_J(\boldsymbol\alpha;t)\}
                          \frac{dt}{\prod_{i\in J}t_i}.
\end{align}
Thus a face contributes only if each of its labels is differentiated
at least once.  The inner jet has order at most
$|\boldsymbol\nu|-|J|+1$.

\begin{corollary}[Squarefree mixed derivatives]
\label{sfj:cor:squarefree}
If $U$ contains $m\ge2$ distinct labels, then
\begin{equation}\label{sfj:eq:squarefree}
 \boxed{\displaystyle
 \left(\prod_{i\in U}\partial_{\alpha_i}\right)
             \widetilde{\mathscr D}_W(0)
 =-\sum_{\substack{J\subseteq U\\|J|\ge2}}
           \frac{I_J(F_{m-|J|+1})}{m-|J|+1}.}
\end{equation}
No logarithmic simplex weights or Gamma Taylor corrections remain.
For three labels this is
\[
 \partial_{\alpha_1}\partial_{\alpha_2}\partial_{\alpha_3}
          \widetilde{\mathscr D}_W(0)
 =-\tfrac12\{I_{12}(F_2)+I_{13}(F_2)+I_{23}(F_2)\}
                                      -I_{123}(F_1).
\]
\end{corollary}
\begin{proof}
In \eqref{sfj:eq:all-mixed}, every variable belonging to $J$ is
already supplied by the face prefactor.  The remaining variables are
outside $J$, whereas $\Pi_J$ contains only variables in $J$.
Consequently it contributes only its constant coefficient.  The
required monomial is supplied by $A^{m-|J|}$ and has coefficient
$(m-|J|)!$.  Division by $(m-|J|+1)!$ gives the stated factor.
\end{proof}

For $\deg W\le2$, this has an ordinary logarithmic-sum interpretation.
Put $L_i(n)=\log((n+a)^2+B_i)$.  With $r\ge2$ labels,
\begin{equation}\label{sfj:eq:ordinary-log-identity}
 \sum_{n=0}^\infty W(n+a)
    \left\{\frac1r\sum_{i=1}^rL_i(n)^r-\prod_{i=1}^rL_i(n)\right\}
 =(-1)^{r+1}
      \sum_{\substack{J\subseteq\{1,\ldots,r\}\\|J|\ge2}}
           \frac{I_J(F_{r-|J|+1})}{r-|J|+1}.
\end{equation}
The leading and first-order changes in $B_i/(n+a)^2$ cancel in the
braces.  The summand is therefore
$O(n^{\deg W-4}\log^{r-2}(2+n))$, and the sum is absolutely
convergent.  This same estimate locally in the order parameters
justifies termwise differentiation of the defect.  The local term is
zero at these polynomial degrees.  Taking the squarefree derivative
in the pointwise defect proves \eqref{sfj:eq:ordinary-log-identity}.

\subsection{Verification and the remaining scope}

The supplied exact implementation compares
\eqref{sfj:eq:monomial} against a separately constructed rational
Dirichlet-moment formula for $2\le k\le5$ and $2\le m\le8$;
all 28 cases agree exactly over the rational numbers.  Six additional
quadratic variance checks cover $2\le k\le7$.

The numerical program uses 48 decimal digits.  It compares direct
triangle face integrals of degrees $3,4,6$ against the finite formula;
their absolute residuals are respectively approximately
$1.02\cdot10^{-26}$, $3.42\cdot10^{-25}$, and
$1.33\cdot10^{-23}$.  The repeated-mass one-dimensional formula and
the independently expanded positive double-polylogarithm series agree
for $k=2,\ldots,6$ to residuals below $1.7\cdot10^{-38}$.
Most significantly, the original \emph{two-dimensional Gamma face}
in \eqref{sfj:eq:numeric-example} was evaluated directly on all six
coordinate-ordering sectors.  At Gauss--Legendre orders $16,24,32$,
its residual against the polylogarithm lattice value decreases from
$2.33\cdot10^{-19}$ to $3.00\cdot10^{-28}$ and then
$4.41\cdot10^{-37}$.  The sector map is
$t=(1,u,uv)/(1+u+uv)$ and its weighted measure is $du\,dv/(uv)$.
The face numerator is retained as a combined expression throughout.

These computations are finite algebraic checks and numerical
diagnostics, not interval certificates or substitutes for the proofs.
The results complete the specified all-order collision and
independent-order extensions relative to the inspected spectral-face
report, and give a global depth-two evaluation in the one-displaced-mass
sector.  They do not prove a finite classical-polylogarithm reduction
for arbitrary distinct masses or every polylogarithm weight.  Natural
next targets are a finite dictionary for two distinct mass clusters,
arithmetic reductions of the resulting double-polylogarithm lattice
sums, and a basis-independent description of the joint polar
subtraction when the affine spectral factors have different degrees.

% END INPUT


% BEGIN sections/08_audit_and_research.tex
\section{Source reconciliation, verification, and further research}
\label{sec:audit}

\subsection{Which earlier questions are settled}

The following table records the precise scope of the continuation.
The source labels refer to the incoming reports at the pinned
snapshot, not to a moving version of the repository.

\begin{longtable}{@{}>{\raggedright\arraybackslash}p{0.27\textwidth}>{\raggedright\arraybackslash}p{0.65\textwidth}@{}}
\toprule
\textbf{Earlier target} & \textbf{Result and remaining scope}\\
\midrule
\endhead
Vertical lattice, Questions 1 and 6 \cite{RepoVertical}
&
Theorems~\ref{thm:minimal-contact} and \ref{thm:filter-all}
give every minimal contact and the exact periodic mean criterion.
All derivative orders are covered.\\[5pt]
Vertical lattice, Question 7 \cite{RepoVertical}
&
Theorem~\ref{thm:modular} gives exact transport for horizontal
shift \(a=1\), near every nonpositive even order and near \(s=1\),
with every derivative correction. A general arbitrary-shift
transport remains open.\\[5pt]
Unequal products beyond two factors \cite{RepoBernoulli,RepoMixed}
&
Theorem~\ref{mtp:generator} handles every finite collection of
positive even tail orders and every regular centered moment.
The separate conjecture about optimal unequal bilinear
prefixes remains open.\\[5pt]
Independent nonperiodic Stieltjes shifts, R8 \cite{RepoTailKernels}
&
Theorem~\ref{gs:completed} supplies an explicit jointly
holomorphic convergent generator, including all argument
derivatives. A finite reduction in a prescribed classical
constant alphabet is a stronger question.\\[5pt]
Higher spectral collisions and independent orders
\cite{RepoBernoulli}
&
The face-moment formula gives all polynomial collision
coefficients. The one-displacement Gamma sector is evaluated
by depth-two polylogarithms. The joint polar part and its
finite local correction are explicit.\\[5pt]
Alternating nonpositive-order square law, R7
\cite{RepoTailKernels}
&
Already resolved in the separate Bilinear Tails archive
\cite{RepoBilinear}; update its status and cite that
proof. It is not counted as a new result here.\\[5pt]
Gaussian \(S_6\) and revised \(S_8\)
\cite{RepoManuscript}
&
Remain conjectural. The present results do not supply the
missing period relations.\\
\bottomrule
\end{longtable}

\subsection{A concrete correction to the research ledger}

In the Tail Kernels archive, R7 asks for the alternating square
continuation at all nonpositive orders. At the same snapshot,
\texttt{sections/04\_cutoff\_jets.tex} in the Bilinear Tails
archive already proves the required normally convergent
continuation, full contact formula, and translated-cutoff rule
\cite{RepoTailKernels,RepoBilinear}. Its theorem labels
\texttt{cut:normal}, \texttt{cut:contacts}, and
\texttt{cut:translation} locate the exact claims.
The order-zero value \(1/4\), despite the vanishing
pointwise summands in the relevant representation, is explicitly
reconciled there. The Gamma specializations have labels
\texttt{cut:Gamma} and \texttt{cut:Gammasquare}.

The proposed replacement ledger entry is:
\begin{quote}
\textbf{R7: resolved.} Cite the all-order cutoff continuation in
\emph{Bilinear Alternating Tails and Harmonic Determinants},
Section 4. Preserve the specified cutoff coordinate and its
order-zero contact. Remaining extensions should be stated
separately from the already established nonpositive-order law.
\end{quote}
This corrects the status shared between two reports. It does not
identify a false theorem in the resolving report.

\subsection{Three safeguards for subsequent integration}

\paragraph{Finite parts must retain their coordinate.}
For the new lattice formula, the constant left after deleting
positive powers of \(\log K\) is
\({\cal F}_M^{(d)}(-N)+C_{N,d}(\log h)\).
Deleting positive powers of the physical-radius logarithm
\(\log R\), \(R=hK\), instead gives
\({\cal F}_M^{(d)}(-N)+C_{N,d}(0)\).
The analytic continuation value is obtained only after the entire
contact polynomial has been removed. These are three specified
operations; their constant differences should not be suppressed
in a shared definition of \emph{regularized sum}.

\paragraph{Periodic cancellation does not imply absolute convergence.}
A nonzero mean-zero periodic multiplier removes the critical
contact at every fixed derivative order. The resulting two
half-lattice series converge by summation by parts. They are in
general conditional, since the terms still have critical
power-log magnitudes. An additional horizontal difference
improves the decay and gives absolute convergence.
Theorem~\ref{thm:filter-all} and
\eqref{eq:relative-lattice} keep these claims distinct.

\paragraph{A common-ray limit is insufficient for independent orders.}
The predecessor's common-ray spectral theorem is correctly
scoped. An unqualified extension to independent orders would
be false. For instance, with \(W(x)=x\), \(a=1\), and masses
\(0,1\), the limits along \((\alpha_1,\alpha_2)=(t,t)\) and
\((t,2t)\) are \(-1/3\) and \(-5/12\).
Even equality of all nontangential straight-ray limits would
not suffice: the example \(W(x)=x^3+x/2\) has polar part
\[
 \frac{\alpha_2^2}{4(\alpha_1+\alpha_2)}.
\]
It tends to zero on every fixed ray with nonzero total direction,
but along \(\alpha_2=t\),
\(\alpha_1=-t+t^3\) it equals \(1/(4t)\).
The joint polar decomposition, rather than a ray calculation,
is the appropriate hypothesis for mixed differentiation.

\subsection{Reproducible verification}

The companion files separate exact algebra from numerical
diagnostics. All arithmetic determinants and finite polynomial
identities are computed over rational or symbolic exact rings.
The floating-point checks compare independently organized
expressions, including direct quadratures, finite cutoff sums,
and Mellin contours.

\begin{longtable}{@{}>{\raggedright\arraybackslash}p{0.25\textwidth}>{\raggedright\arraybackslash}p{0.67\textwidth}@{}}
\toprule
\textbf{Component} & \textbf{Verification supplied}\\
\midrule
\endhead
Lattice contacts
&
60 exact coefficient and Pochhammer checks, two mutation
controls, and 27 finite-lattice diagnostics at 55 digits.
The next two Euler--Maclaurin coefficients are used to check
the remainder when its leading polynomial nearly cancels.\\[4pt]
Modular transport
&
Nine exact kernel-polynomial checks, four direct
Lambert-versus-contour comparisons at 42 digits,
two classical boundary checks, and an independent
self-dual contour evaluation.\\[4pt]
Mixed tails
&
306 exact moments for all multisets of one through four
orders from \(\{2,4,6\}\), through the tested moment range,
against finite summation by parts; 825 boundary-parity
checks; exact six- and eight-coordinate determinants;
a finite telescoping certificate; and 90-digit checks
of the short identities.\\[4pt]
Independent shifts
&
16 exact cancellation checks and 42 numerical checks
at 65 digits, including complex shifts, both spectral
axes, \(u+v=0\), prefix invariance, and the explicit
\(\gamma_1\times\psi\) coefficient.
The largest recorded absolute residual is approximately
\(1.02\times10^{-56}\).\\[4pt]
Spectral faces
&
34 exact rational checks and 11 numerical comparisons
at 48-digit working precision. The original
two-dimensional Gamma face is checked against the
polylogarithmic lattice expression with increasing
quadrature orders.\\
\bottomrule
\end{longtable}

These calculations are regression diagnostics for the stated
proofs. They do not establish global arithmetic independence,
replace a normal-convergence argument, or certify decimal
enclosures. All domains, branch conventions, and limiting
operations used in the identities are established in the text.
The verification programs and their recorded outputs are
included so that future integration can detect sign, endpoint,
normalization, or coefficient regressions.

\subsection{Further research questions}
\label{sec:questions}

The following questions are ordered by their connection to the
new mechanisms and by the prospect of obtaining further exact
identities. Each asks for a structural statement or evaluation.

\begin{enumerate}[label=\textbf{Q\arabic*.},leftmargin=*]
\item \textbf{Optimal prefixes for arbitrary unequal tails.}
The mixed recurrence proves a finite span, and the two worked
products have complete rational prefix certificates.
Determine the minimal formal coordinate dimension for a
general even-order multiset \(I\), after a declared system of
MZV reductions, and determine when its first admissible
moments already span that space. The existing unequal
bilinear optimal-prefix conjecture is a particularly precise
first target. A proof should establish an explicit nonzero
determinant, rather than infer rank from evaluated decimals.

\item \textbf{Direct telescoping certificates for elegant tail identities.}
The short cubic identity admits a finite telescoping
certificate in products of Hurwitz tails. Characterize the
centered polynomial combinations for which such a certificate
exists. Can the polarized recurrence be compiled into a
finite certificate with a prescribed maximal tail order?
This could systematically produce short ordinary identities
whose cancellations hide a much larger moment system.

\item \textbf{Minimal functional depth of mixed generators.}
The bound \(d-1\) is constructive but need not be sharp.
Classify the order multisets for which the anchored generator
reduces to classical polylogarithms, or to depth two
independently of \(d\). The classification must distinguish
functional identities from numerical MZV reductions at the
endpoint \(t=0\).

\item \textbf{More than one displaced spectral mass.}
For \(k-1\) equal masses and one displaced mass, depth two
suffices in every dimension. Determine the exact function
class for two displaced masses and, more generally, for a
fixed number of distinct mass clusters. Is depth controlled
by the number of clusters rather than the number of factors?
A useful result would give a finite multiple-polylogarithm
formula with explicit branches and ordinary convergence,
not merely rename the remaining simplex integral.

\item \textbf{Complete logarithmic collision coefficients.}
The rational normalized face formula produces every
logarithmically weighted polynomial moment by differentiation.
Find compact closed forms for the first genuinely nonsquarefree
mixed order derivatives, and organize their coefficients by
harmonic star sums and zeta weight. Determine which Gamma
normalization coefficients cancel before any MZV reduction.

\item \textbf{Independent-order continuation beyond quadratic spectra.}
For quadratic factors, the odd powers of \(W\) give the exact
polar obstruction through a single total-order divisor.
Develop the corresponding finite polar-part classification
for higher-degree polynomial spectra and shifted linear
factors. Identify the natural local subtraction that is
compatible with multiplication, parameter differentiation,
and periodic mean cancellation.

\item \textbf{Modular transport with arbitrary horizontal shift.}
Extend the exact order-dependent contour identity from
\(a=1\) to general \(a\). Rational shifts should first be
organized into finite systems with additive characters and
Hurwitz factors. For irrational shifts, determine the
additional integral transform explicitly. The fixed-cutoff
contact polynomial should remain visible throughout the
transport.

\item \textbf{Arithmetic reductions of the logarithmic Hurwitz family.}
For which rational pairs \((c,z)\) does
\[
 \mathcal L_p(c,z)=
 \sum_{n\ge0}\frac{\log(1+z/(n+c))}{(n+c)^p}
\]
reduce to a finite specified alphabet of polylogarithms,
Gamma values, and zeta derivatives?
The generating and differential relations in the
independent-shift section impose strong compatibility
conditions. Generic rational arguments should not be
assumed to have such a reduction.

\item \textbf{Unequal-rate Stieltjes boundary collisions.}
Set \(a=\lambda\epsilon\), \(b=\mu\epsilon\) with fixed
\(\lambda,\mu>0\), and let \(\epsilon\downarrow0\).
Compute the complete singular power-log primitive and
the remaining constant for arbitrary Stieltjes indices
and argument derivatives. Determine the exact path
correction relative to the common-shift finite part.
The diagonal unfolding and explicit finite-prefix
derivatives provide a starting point.

\item \textbf{Three independently shifted Stieltjes factors.}
Extend the evaluated two-factor germ to three factors
with independent shifts and orders. Partitioning the
summation lattice into minimum-index regions suggests
finite Lauricella-type prefixes. The central problem is
to remove all intersecting spectral poles jointly while
retaining geometric convergence after every derivative.

\item \textbf{The unresolved Gaussian period identities.}
The manuscript's \(S_6\) and revised \(S_8\) candidates
need exact period relations. The existing high-precision
residuals and bounded relation-space calculations do
not supply them. Seek a new functional relation or
contour deformation that escapes the already tested
relation spaces, and exhibit its exact coefficient
certificate before numerical evaluation. The new
mixed-tail rank results concern different spaces and
should not be used as evidence of a Gaussian proof.

\item \textbf{A formal proof layer for the finite core.}
The nilpotent polarized Bernoulli identity, the exact
boundary-parity cancellation, the contact-polynomial
coefficient identity, and the rational face compiler
are natural Lean targets. Separate their finite
algebraic statements from Euler--Maclaurin continuation
and normal-convergence lemmas. This would make later
analytic generalizations depend on a small reusable
verified algebraic core, without claiming that the
current Python checks constitute formal proofs.
\end{enumerate}


% END INPUT

\clearpage
\begin{thebibliography}{99}

% BEGIN references.tex
\bibitem{RepoManuscript}
Vladimir Reshetnikov's ProveIt project,
\emph{Analysis/Polylogarithms: manuscript, editorial ledger, and source chapters}.
Snapshot \texttt{36edefa8c531d9d6df8306ac51bcc317250f404c},
11 October 2026.
\href{https://github.com/VladimirReshetnikov/ProveIt/tree/36edefa8c531d9d6df8306ac51bcc317250f404c/Analysis/Polylogarithms/docs/manuscript}{Pinned manuscript directory}.
The Gaussian status is recorded in its README and
\texttt{chapters/04-S8-candidate.tex}.

\bibitem{RepoIncoming}
ProveIt project, \emph{Incoming research reports},
the ten archives present in the same pinned snapshot.
\href{https://github.com/VladimirReshetnikov/ProveIt/tree/36edefa8c531d9d6df8306ac51bcc317250f404c/docs/incoming}{Pinned incoming directory}.
Exact archive names, Git blob identifiers, and SHA-256 digests are
recorded in the companion provenance file.

\bibitem{RepoVertical}
ProveIt research continuation,
\emph{Discrete Vertical Hurwitz Calculus:
Stieltjes jets, polylogarithmic sampling, and resonant contact terms},
11 October 2026; especially the compensated sampling theorem and
research questions 1, 6, and 7.
\href{https://github.com/VladimirReshetnikov/ProveIt/blob/36edefa8c531d9d6df8306ac51bcc317250f404c/docs/incoming/ProveIt_Vertical_Lattice_Jets_2026-10-11.zip}{Pinned Vertical Lattice Jets archive}.

\bibitem{RepoBernoulli}
ProveIt research continuation,
\emph{Bernoulli Cancellation, Harmonic Gaps, and Spectral Faces},
11 October 2026. Section 2 gives the repeated-tail recurrence and
explicitly separates unequal products of three or more factors;
Section 5 gives the common-ray face formula and further targets.
\href{https://github.com/VladimirReshetnikov/ProveIt/blob/36edefa8c531d9d6df8306ac51bcc317250f404c/docs/incoming/ProveIt_Bernoulli_Tails_and_Spectral_Faces_2026-10-11.zip}{Pinned Bernoulli Tails and Spectral Faces archive}.

\bibitem{RepoMixed}
ProveIt research continuation,
\emph{Harmonic Gaps and Mixed Moments}, Section 4,
\emph{All mixed even-tail moments}.
Source in the harmonic-gaps-mixed-moments continuation,
at the same pinned snapshot.
\href{https://github.com/VladimirReshetnikov/ProveIt/blob/36edefa8c531d9d6df8306ac51bcc317250f404c/Analysis/Polylogarithms/docs/reports/stieltjes-harmonic-resonance/continuations/harmonic-gaps-mixed-moments/sections/04_mixed_moments.tex}{Pinned mixed-moment source}.

\bibitem{RepoTailKernels}
ProveIt research continuation,
\emph{Bilinear Tail Identities, Polylogarithmic Resonances,
and Stieltjes Boundary Products},
11 October 2026; especially research questions R7--R8.
\href{https://github.com/VladimirReshetnikov/ProveIt/blob/36edefa8c531d9d6df8306ac51bcc317250f404c/docs/incoming/ProveIt_Tail_Kernels_and_Boundary_Products_2026-10-11.zip}{Pinned Tail Kernels and Boundary Products archive}.

\bibitem{RepoBilinear}
ProveIt research continuation,
\emph{Bilinear Alternating Tails and Harmonic Determinants},
11 October 2026; especially
\texttt{sections/04\_cutoff\_jets.tex}.
\href{https://github.com/VladimirReshetnikov/ProveIt/blob/36edefa8c531d9d6df8306ac51bcc317250f404c/docs/incoming/ProveIt_Bilinear_Tails_and_Determinants_2026-10-11.zip}{Pinned Bilinear Tails and Determinants archive}.

\bibitem{DLMF}
F.~W.~J. Olver and others, eds.,
\emph{NIST Digital Library of Mathematical Functions}.
Sections 5.7--5.8 (Gamma expansions and products), 5.15 (polygamma),
24.4 (Bernoulli identities), 25.4 (zeta functional equation),
25.11 (Hurwitz zeta), and 25.14 (Lerch transcendent).
\url{https://dlmf.nist.gov/}.
Accessed 11 October 2026.

\bibitem{Hoffman}
M.~E. Hoffman,
\emph{Sums of products of Riemann zeta tails},
Mediterranean Journal of Mathematics \textbf{13} (2016), 2771--2781.
\url{https://doi.org/10.1007/s00009-015-0653-9}.
\url{https://arxiv.org/abs/1602.03399}.

\bibitem{LiChu}
C. Li and W. Chu,
\emph{Product sums of zeta tails by series rearrangements},
Computational Mathematics and Mathematical Physics
\textbf{65} (2025), 8--24.
\url{https://doi.org/10.1134/S0965542524701896}.
The publisher's abstract and bibliographical record were consulted;
no claim about the uninspected full text is used in a proof.

\bibitem{ShpotParis}
M.~A. Shpot and R.~B. Paris,
\emph{Integrals of products of Hurwitz zeta functions via Feynman
parametrization and two double sums of Riemann zeta functions},
arXiv:1609.05658v2, 24 May 2020.
\url{https://arxiv.org/abs/1609.05658}.

\bibitem{Connon}
D.~F. Connon,
\emph{On an integral involving the digamma function},
arXiv:1212.1432v2, 9 December 2012.
\url{https://arxiv.org/abs/1212.1432}.

\bibitem{Carlson}
B.~C. Carlson,
\emph{Dirichlet averages of \(x^t\log x\)},
SIAM Journal on Mathematical Analysis \textbf{18} (1987), 550--565.
\url{https://doi.org/10.1137/0518043}.
The primary publisher abstract and bibliographical record were
inspected; this citation credits the classical Dirichlet-average
context rather than claiming a full-paper priority audit.

\bibitem{ArdehaliRosengren}
A. Arabi Ardehali and H. Rosengren,
\emph{A new product formula for \((z;q)_\infty\),
with applications to asymptotics},
arXiv:2602.11329v2, 26 February 2026.
\url{https://arxiv.org/abs/2602.11329}.
Published version:
\url{https://doi.org/10.1007/s00365-026-09783-2}.

\bibitem{DixitStieltjes}
A. Dixit, S. Sathyanarayana, and N. Guru Sharan,
\emph{Modular relations involving generalized digamma functions},
arXiv:2306.10991v2, 22 June 2023;
Journal of Mathematical Analysis and Applications
\textbf{539} (2024), article 128479.
\url{https://arxiv.org/abs/2306.10991}.


% END INPUT

\end{thebibliography}
\end{document}

